% ---- score2 standard preamble (from: python3 score2_tools.py preamble) ----
\documentclass[10pt]{extarticle}
\usepackage[a4paper,margin=1.5cm]{geometry}
\usepackage[T1]{fontenc}
\usepackage{amsmath,amssymb}
\usepackage{graphicx}
\usepackage[export]{adjustbox}
\usepackage{multicol}
\usepackage{enumitem}
\usepackage{booktabs,array,tabularx}
\setlength{\columnsep}{0.7cm}
\setlength{\parindent}{0pt}
\setlength{\parskip}{3pt}
\pagestyle{empty}
\makeatletter
% figure registry (figdefs.tex is generated by the extract step)
\newcommand{\qfigdef}[4]{\@namedef{qf@file@#1}{#2}\@namedef{qf@w@#1}{#3}\@namedef{qf@wide@#1}{#4}}
\newcommand{\qf@check}[1]{\@ifundefined{qf@file@#1}{\PackageError{score2}{Unknown figure id #1}{}}{}}
\newcommand{\qf@img}[4][]{\includegraphics[width=#2\dimexpr\@nameuse{qf@w@#4}pt\relax,max width=#3,max height=0.45\textheight,#1]{\@nameuse{qf@file@#4}}}
% Full-width questions: finalize wraps every question that contains a wide figure in \qfullbegin...\qfullend,
% so the WHOLE question (stem, figures, options) is set across the page in reading order, then the two columns
% resume. The question's head - its number, stem and first figure (\qfullsplit marks the end of the head) -
% is kept on one page, so a number is never left at a page foot with its figure overleaf; the rest may break
% across pages normally, so no big gaps are left. Without a \qfullsplit, the whole question is kept together
% if it is under 45% of a page. Never write these yourself.
\newif\ifqf@full
\newif\ifqf@boxed
\newsavebox{\qf@box}
\newcommand{\qf@open}{\global\qf@fulltrue\par\global\qf@boxedtrue
  \setbox\qf@box\vbox\bgroup\hsize\textwidth\linewidth\textwidth}
\newcommand{\qf@emit}[1]{\par\egroup\global\qf@boxedfalse
  \ifdim\dimexpr\ht\qf@box+\dp\qf@box\relax<#1\textheight\noindent\usebox{\qf@box}\par\else\unvbox\qf@box\fi}
\newcommand{\qfullsplit}{\ifqf@boxed\qf@emit{0.85}\fi}
\newcommand{\qf@close}{\ifqf@boxed\qf@emit{0.45}\fi\global\qf@fullfalse}
\newcommand{\qfullbegin}{\end{multicols}\qf@open}
\newcommand{\qfullend}{\qf@close\begin{multicols}{2}}
\newcommand{\qfullstart}{\qf@open}
\newcommand{\qfulldone}{\qf@close}
% \qfig[scale]{id}: block figure, centred, directly where it is written (size set by the tool).
% Fallback if a wide figure is not inside a full-width question: break out of the columns in place - no floats.
\newcommand{\qfig}[2][1]{\qf@check{#2}%
  \ifqf@full\qf@blk{#1}{#2}\else
  \ifnum\@nameuse{qf@wide@#2}=1 \expandafter\@firstoftwo\else\expandafter\@secondoftwo\fi
  {\end{multicols}\begin{center}\qf@img{#1}{\textwidth}{#2}\end{center}\begin{multicols}{2}}%
  {\qf@blk{#1}{#2}}\fi}
\newcommand{\qf@blk}[2]{\par\smallskip{\centering\qf@img{#1}{\linewidth}{#2}\par}\smallskip}
% \qfigi[scale]{id}: inline figure (inside option lists / tables / minipages); never breaks columns.
\newcommand{\qfigi}[2][1]{\qf@check{#2}\qf@img[valign=c]{#1}{\linewidth}{#2}}
\makeatother
% structure macros - lint reads these to check completeness and figure placement
\newcommand{\papersec}[1]{\par\bigskip{\centering\Large\bfseries PAPER-#1\par}\medskip}
\newcommand{\partsec}[2]{\par\medskip\noindent{\large\bfseries PART-#1: #2}\par\smallskip}
\newcommand{\sect}[1]{\par\smallskip\noindent\textbf{#1}\par}
\newcommand{\qn}[1]{\par\smallskip\noindent\textbf{#1.}~}
\newcommand{\paraq}[2]{\par\smallskip\noindent\textbf{Paragraph for Questions #1 and #2}\par}
\newenvironment{qoptions}{\begin{itemize}[leftmargin=1.8em,itemsep=1pt,topsep=2pt,label={}]}{\end{itemize}}
\newcommand{\opt}[2]{\item[(#1)] #2}
% figure registry (generated by score2_tools.py finalize)
\qfigdef{p1_phy_s1_q01_1}{Pre-Score (2023-24) - Test 03 (22 Oct 2023)-img/p1_phy_s1_q01_1.png}{65.8}{0}
\qfigdef{p1_phy_s1_q03_1}{Pre-Score (2023-24) - Test 03 (22 Oct 2023)-img/p1_phy_s1_q03_1.png}{200.8}{0}
\qfigdef{p1_phy_s2_q08_1}{Pre-Score (2023-24) - Test 03 (22 Oct 2023)-img/p1_phy_s2_q08_1.png}{131.5}{0}
\qfigdef{p1_phy_s2_q09_1}{Pre-Score (2023-24) - Test 03 (22 Oct 2023)-img/p1_phy_s2_q09_1.png}{157.9}{0}
\qfigdef{p1_phy_s2_q10_1}{Pre-Score (2023-24) - Test 03 (22 Oct 2023)-img/p1_phy_s2_q10_1.png}{63.5}{0}
\qfigdef{p1_phy_s2_q12_1}{Pre-Score (2023-24) - Test 03 (22 Oct 2023)-img/p1_phy_s2_q12_1.png}{103.1}{0}
\qfigdef{p1_phy_s3_q02_1}{Pre-Score (2023-24) - Test 03 (22 Oct 2023)-img/p1_phy_s3_q02_1.png}{97.0}{0}
\qfigdef{p1_phy_s3_q03_1}{Pre-Score (2023-24) - Test 03 (22 Oct 2023)-img/p1_phy_s3_q03_1.png}{102.3}{0}
\qfigdef{p1_phy_s3_q04_1}{Pre-Score (2023-24) - Test 03 (22 Oct 2023)-img/p1_phy_s3_q04_1.png}{149.7}{0}
\qfigdef{p1_phy_s3_q04_2}{Pre-Score (2023-24) - Test 03 (22 Oct 2023)-img/p1_phy_s3_q04_2.png}{161.1}{0}
\qfigdef{p1_phy_s3_q05_1}{Pre-Score (2023-24) - Test 03 (22 Oct 2023)-img/p1_phy_s3_q05_1.png}{111.4}{0}
\qfigdef{p1_chem_s1_q01_1}{Pre-Score (2023-24) - Test 03 (22 Oct 2023)-img/p1_chem_s1_q01_1.png}{192.3}{0}
\qfigdef{p1_chem_s1_q01_2}{Pre-Score (2023-24) - Test 03 (22 Oct 2023)-img/p1_chem_s1_q01_2.png}{117.5}{0}
\qfigdef{p1_chem_s1_q01_A}{Pre-Score (2023-24) - Test 03 (22 Oct 2023)-img/p1_chem_s1_q01_A.png}{62.9}{0}
\qfigdef{p1_chem_s1_q01_B}{Pre-Score (2023-24) - Test 03 (22 Oct 2023)-img/p1_chem_s1_q01_B.png}{55.3}{0}
\qfigdef{p1_chem_s1_q01_C}{Pre-Score (2023-24) - Test 03 (22 Oct 2023)-img/p1_chem_s1_q01_C.png}{65.1}{0}
\qfigdef{p1_chem_s1_q01_D}{Pre-Score (2023-24) - Test 03 (22 Oct 2023)-img/p1_chem_s1_q01_D.png}{71.7}{0}
\qfigdef{p1_chem_s1_q02_1}{Pre-Score (2023-24) - Test 03 (22 Oct 2023)-img/p1_chem_s1_q02_1.png}{160.7}{0}
\qfigdef{p1_chem_s1_q03_1}{Pre-Score (2023-24) - Test 03 (22 Oct 2023)-img/p1_chem_s1_q03_1.png}{85.7}{0}
\qfigdef{p1_chem_s1_q03_A}{Pre-Score (2023-24) - Test 03 (22 Oct 2023)-img/p1_chem_s1_q03_A.png}{142.5}{0}
\qfigdef{p1_chem_s1_q03_B}{Pre-Score (2023-24) - Test 03 (22 Oct 2023)-img/p1_chem_s1_q03_B.png}{138.7}{0}
\qfigdef{p1_chem_s1_q03_C}{Pre-Score (2023-24) - Test 03 (22 Oct 2023)-img/p1_chem_s1_q03_C.png}{141.0}{0}
\qfigdef{p1_chem_s1_q03_D}{Pre-Score (2023-24) - Test 03 (22 Oct 2023)-img/p1_chem_s1_q03_D.png}{108.2}{0}
\qfigdef{p1_chem_s1_q04_1}{Pre-Score (2023-24) - Test 03 (22 Oct 2023)-img/p1_chem_s1_q04_1.png}{245.2}{0}
\qfigdef{p1_chem_s1_q04_2}{Pre-Score (2023-24) - Test 03 (22 Oct 2023)-img/p1_chem_s1_q04_2.png}{267.2}{1}
\qfigdef{p1_chem_s2_q05_A}{Pre-Score (2023-24) - Test 03 (22 Oct 2023)-img/p1_chem_s2_q05_A.png}{122.9}{0}
\qfigdef{p1_chem_s2_q05_B}{Pre-Score (2023-24) - Test 03 (22 Oct 2023)-img/p1_chem_s2_q05_B.png}{164.3}{0}
\qfigdef{p1_chem_s2_q05_C}{Pre-Score (2023-24) - Test 03 (22 Oct 2023)-img/p1_chem_s2_q05_C.png}{193.4}{0}
\qfigdef{p1_chem_s2_q05_D}{Pre-Score (2023-24) - Test 03 (22 Oct 2023)-img/p1_chem_s2_q05_D.png}{95.1}{0}
\qfigdef{p1_chem_s2_q06_1}{Pre-Score (2023-24) - Test 03 (22 Oct 2023)-img/p1_chem_s2_q06_1.png}{115.1}{0}
\qfigdef{p1_chem_s2_q06_B}{Pre-Score (2023-24) - Test 03 (22 Oct 2023)-img/p1_chem_s2_q06_B.png}{41.4}{0}
\qfigdef{p1_chem_s2_q06_A}{Pre-Score (2023-24) - Test 03 (22 Oct 2023)-img/p1_chem_s2_q06_A.png}{30.3}{0}
\qfigdef{p1_chem_s2_q06_C}{Pre-Score (2023-24) - Test 03 (22 Oct 2023)-img/p1_chem_s2_q06_C.png}{50.0}{0}
\qfigdef{p1_chem_s2_q06_D}{Pre-Score (2023-24) - Test 03 (22 Oct 2023)-img/p1_chem_s2_q06_D.png}{35.4}{0}
\qfigdef{p1_chem_s2_q07_A}{Pre-Score (2023-24) - Test 03 (22 Oct 2023)-img/p1_chem_s2_q07_A.png}{166.5}{0}
\qfigdef{p1_chem_s2_q07_B}{Pre-Score (2023-24) - Test 03 (22 Oct 2023)-img/p1_chem_s2_q07_B.png}{200.1}{0}
\qfigdef{p1_chem_s2_q07_C}{Pre-Score (2023-24) - Test 03 (22 Oct 2023)-img/p1_chem_s2_q07_C.png}{98.0}{0}
\qfigdef{p1_chem_s2_q07_4}{Pre-Score (2023-24) - Test 03 (22 Oct 2023)-img/p1_chem_s2_q07_4.png}{55.9}{0}
\qfigdef{p1_chem_s2_q07_D}{Pre-Score (2023-24) - Test 03 (22 Oct 2023)-img/p1_chem_s2_q07_D.png}{216.8}{0}
\qfigdef{p1_chem_s2_q08_1}{Pre-Score (2023-24) - Test 03 (22 Oct 2023)-img/p1_chem_s2_q08_1.png}{144.9}{0}
\qfigdef{p1_chem_s2_q08_2}{Pre-Score (2023-24) - Test 03 (22 Oct 2023)-img/p1_chem_s2_q08_2.png}{23.8}{0}
\qfigdef{p1_chem_s2_q08_3}{Pre-Score (2023-24) - Test 03 (22 Oct 2023)-img/p1_chem_s2_q08_3.png}{23.8}{0}
\qfigdef{p1_chem_s2_q09_A}{Pre-Score (2023-24) - Test 03 (22 Oct 2023)-img/p1_chem_s2_q09_A.png}{27.0}{0}
\qfigdef{p1_chem_s2_q09_B}{Pre-Score (2023-24) - Test 03 (22 Oct 2023)-img/p1_chem_s2_q09_B.png}{51.5}{0}
\qfigdef{p1_chem_s2_q09_C}{Pre-Score (2023-24) - Test 03 (22 Oct 2023)-img/p1_chem_s2_q09_C.png}{26.5}{0}
\qfigdef{p1_chem_s2_q09_D}{Pre-Score (2023-24) - Test 03 (22 Oct 2023)-img/p1_chem_s2_q09_D.png}{59.5}{0}
\qfigdef{p1_chem_s2_q10_1}{Pre-Score (2023-24) - Test 03 (22 Oct 2023)-img/p1_chem_s2_q10_1.png}{37.3}{0}
\qfigdef{p1_chem_s2_q10_A}{Pre-Score (2023-24) - Test 03 (22 Oct 2023)-img/p1_chem_s2_q10_A.png}{58.2}{0}
\qfigdef{p1_chem_s2_q10_C}{Pre-Score (2023-24) - Test 03 (22 Oct 2023)-img/p1_chem_s2_q10_C.png}{98.5}{0}
\qfigdef{p1_chem_s2_q10_D}{Pre-Score (2023-24) - Test 03 (22 Oct 2023)-img/p1_chem_s2_q10_D.png}{113.9}{0}
\qfigdef{p1_chem_s2_q10_5}{Pre-Score (2023-24) - Test 03 (22 Oct 2023)-img/p1_chem_s2_q10_5.png}{48.5}{0}
\qfigdef{p1_chem_s2_q11_1}{Pre-Score (2023-24) - Test 03 (22 Oct 2023)-img/p1_chem_s2_q11_1.png}{170.8}{0}
\qfigdef{p1_chem_s2_q11_2}{Pre-Score (2023-24) - Test 03 (22 Oct 2023)-img/p1_chem_s2_q11_2.png}{226.1}{0}
\qfigdef{p1_chem_s2_q12_1}{Pre-Score (2023-24) - Test 03 (22 Oct 2023)-img/p1_chem_s2_q12_1.png}{167.5}{0}
\qfigdef{p1_chem_s2_q12_2}{Pre-Score (2023-24) - Test 03 (22 Oct 2023)-img/p1_chem_s2_q12_2.png}{161.8}{0}
\qfigdef{p1_chem_s2_q12_C}{Pre-Score (2023-24) - Test 03 (22 Oct 2023)-img/p1_chem_s2_q12_C.png}{114.4}{0}
\qfigdef{p1_chem_s2_q12_D}{Pre-Score (2023-24) - Test 03 (22 Oct 2023)-img/p1_chem_s2_q12_D.png}{129.3}{0}
\qfigdef{p1_chem_s3_q01_1}{Pre-Score (2023-24) - Test 03 (22 Oct 2023)-img/p1_chem_s3_q01_1.png}{181.2}{0}
\qfigdef{p1_chem_s3_q01_2}{Pre-Score (2023-24) - Test 03 (22 Oct 2023)-img/p1_chem_s3_q01_2.png}{155.8}{0}
\qfigdef{p1_chem_s3_q01_3}{Pre-Score (2023-24) - Test 03 (22 Oct 2023)-img/p1_chem_s3_q01_3.png}{174.6}{0}
\qfigdef{p1_chem_s3_q02_1}{Pre-Score (2023-24) - Test 03 (22 Oct 2023)-img/p1_chem_s3_q02_1.png}{266.7}{1}
\qfigdef{p1_chem_s3_q02_3}{Pre-Score (2023-24) - Test 03 (22 Oct 2023)-img/p1_chem_s3_q02_3.png}{215.4}{0}
\qfigdef{p1_chem_s3_q02_4}{Pre-Score (2023-24) - Test 03 (22 Oct 2023)-img/p1_chem_s3_q02_4.png}{272.3}{1}
\qfigdef{p1_chem_s3_q02_5}{Pre-Score (2023-24) - Test 03 (22 Oct 2023)-img/p1_chem_s3_q02_5.png}{273.3}{1}
\qfigdef{p1_chem_s3_q03_1}{Pre-Score (2023-24) - Test 03 (22 Oct 2023)-img/p1_chem_s3_q03_1.png}{315.9}{1}
\qfigdef{p1_chem_s3_q03_2}{Pre-Score (2023-24) - Test 03 (22 Oct 2023)-img/p1_chem_s3_q03_2.png}{285.6}{1}
\qfigdef{p1_chem_s3_q03_3}{Pre-Score (2023-24) - Test 03 (22 Oct 2023)-img/p1_chem_s3_q03_3.png}{206.2}{0}
\qfigdef{p1_chem_s3_q03_4}{Pre-Score (2023-24) - Test 03 (22 Oct 2023)-img/p1_chem_s3_q03_4.png}{159.7}{0}
\qfigdef{p1_chem_s3_q04_1}{Pre-Score (2023-24) - Test 03 (22 Oct 2023)-img/p1_chem_s3_q04_1.png}{215.0}{0}
\qfigdef{p1_chem_s3_q04_2}{Pre-Score (2023-24) - Test 03 (22 Oct 2023)-img/p1_chem_s3_q04_2.png}{223.8}{0}
\qfigdef{p1_chem_s3_q04_3}{Pre-Score (2023-24) - Test 03 (22 Oct 2023)-img/p1_chem_s3_q04_3.png}{245.2}{0}
\qfigdef{p1_chem_s3_q05_1}{Pre-Score (2023-24) - Test 03 (22 Oct 2023)-img/p1_chem_s3_q05_1.png}{42.4}{0}
\qfigdef{p1_chem_s3_q06_1}{Pre-Score (2023-24) - Test 03 (22 Oct 2023)-img/p1_chem_s3_q06_1.png}{130.0}{0}
\qfigdef{p1_chem_s3_q06_2}{Pre-Score (2023-24) - Test 03 (22 Oct 2023)-img/p1_chem_s3_q06_2.png}{186.9}{0}
\qfigdef{p1_chem_s3_q06_3}{Pre-Score (2023-24) - Test 03 (22 Oct 2023)-img/p1_chem_s3_q06_3.png}{29.4}{0}
\qfigdef{p1_chem_s3_q06_4}{Pre-Score (2023-24) - Test 03 (22 Oct 2023)-img/p1_chem_s3_q06_4.png}{156.0}{0}
\qfigdef{p2_phy_s1_q01_1}{Pre-Score (2023-24) - Test 03 (22 Oct 2023)-img/p2_phy_s1_q01_1.png}{186.6}{0}
\qfigdef{p2_phy_s1_q02_1}{Pre-Score (2023-24) - Test 03 (22 Oct 2023)-img/p2_phy_s1_q02_1.png}{108.2}{0}
\qfigdef{p2_phy_s1_q04_1}{Pre-Score (2023-24) - Test 03 (22 Oct 2023)-img/p2_phy_s1_q04_1.png}{140.1}{0}
\qfigdef{p2_phy_s1_q05_A}{Pre-Score (2023-24) - Test 03 (22 Oct 2023)-img/p2_phy_s1_q05_A.png}{76.4}{0}
\qfigdef{p2_phy_s1_q05_B}{Pre-Score (2023-24) - Test 03 (22 Oct 2023)-img/p2_phy_s1_q05_B.png}{76.4}{0}
\qfigdef{p2_phy_s1_q05_C}{Pre-Score (2023-24) - Test 03 (22 Oct 2023)-img/p2_phy_s1_q05_C.png}{76.1}{0}
\qfigdef{p2_phy_s1_q05_D}{Pre-Score (2023-24) - Test 03 (22 Oct 2023)-img/p2_phy_s1_q05_D.png}{76.1}{0}
\qfigdef{p2_phy_s1_q06_1}{Pre-Score (2023-24) - Test 03 (22 Oct 2023)-img/p2_phy_s1_q06_1.png}{213.9}{0}
\qfigdef{p2_phy_s1_q08_1}{Pre-Score (2023-24) - Test 03 (22 Oct 2023)-img/p2_phy_s1_q08_1.png}{95.1}{0}
\qfigdef{p2_phy_s2_q09_1}{Pre-Score (2023-24) - Test 03 (22 Oct 2023)-img/p2_phy_s2_q09_1.png}{107.5}{0}
\qfigdef{p2_phy_s2_q10_1}{Pre-Score (2023-24) - Test 03 (22 Oct 2023)-img/p2_phy_s2_q10_1.png}{107.5}{0}
\qfigdef{p2_phy_s2_q11_1}{Pre-Score (2023-24) - Test 03 (22 Oct 2023)-img/p2_phy_s2_q11_1.png}{65.9}{0}
\qfigdef{p2_phy_s2_q12_1}{Pre-Score (2023-24) - Test 03 (22 Oct 2023)-img/p2_phy_s2_q12_1.png}{63.2}{0}
\qfigdef{p2_phy_s3_q02_1}{Pre-Score (2023-24) - Test 03 (22 Oct 2023)-img/p2_phy_s3_q02_1.png}{149.9}{0}
\qfigdef{p2_phy_s3_q05_1}{Pre-Score (2023-24) - Test 03 (22 Oct 2023)-img/p2_phy_s3_q05_1.png}{311.1}{1}
\qfigdef{p2_phy_s3_q06_1}{Pre-Score (2023-24) - Test 03 (22 Oct 2023)-img/p2_phy_s3_q06_1.png}{121.5}{0}
\qfigdef{p2_chem_s1_q01_1}{Pre-Score (2023-24) - Test 03 (22 Oct 2023)-img/p2_chem_s1_q01_1.png}{76.5}{0}
\qfigdef{p2_chem_s1_q01_2}{Pre-Score (2023-24) - Test 03 (22 Oct 2023)-img/p2_chem_s1_q01_2.png}{136.5}{0}
\qfigdef{p2_chem_s1_q02_A}{Pre-Score (2023-24) - Test 03 (22 Oct 2023)-img/p2_chem_s1_q02_A.png}{47.7}{0}
\qfigdef{p2_chem_s1_q02_B}{Pre-Score (2023-24) - Test 03 (22 Oct 2023)-img/p2_chem_s1_q02_B.png}{32.3}{0}
\qfigdef{p2_chem_s1_q02_C}{Pre-Score (2023-24) - Test 03 (22 Oct 2023)-img/p2_chem_s1_q02_C.png}{122.2}{0}
\qfigdef{p2_chem_s1_q02_D}{Pre-Score (2023-24) - Test 03 (22 Oct 2023)-img/p2_chem_s1_q02_D.png}{40.9}{0}
\qfigdef{p2_chem_s1_q05_1}{Pre-Score (2023-24) - Test 03 (22 Oct 2023)-img/p2_chem_s1_q05_1.png}{110.7}{0}
\qfigdef{p2_chem_s1_q06_A}{Pre-Score (2023-24) - Test 03 (22 Oct 2023)-img/p2_chem_s1_q06_A.png}{101.8}{0}
\qfigdef{p2_chem_s1_q06_B}{Pre-Score (2023-24) - Test 03 (22 Oct 2023)-img/p2_chem_s1_q06_B.png}{98.0}{0}
\qfigdef{p2_chem_s1_q06_C}{Pre-Score (2023-24) - Test 03 (22 Oct 2023)-img/p2_chem_s1_q06_C.png}{92.8}{0}
\qfigdef{p2_chem_s1_q06_D}{Pre-Score (2023-24) - Test 03 (22 Oct 2023)-img/p2_chem_s1_q06_D.png}{95.8}{0}
\qfigdef{p2_chem_s1_q07_1}{Pre-Score (2023-24) - Test 03 (22 Oct 2023)-img/p2_chem_s1_q07_1.png}{136.9}{0}
\qfigdef{p2_chem_s2_q09_1}{Pre-Score (2023-24) - Test 03 (22 Oct 2023)-img/p2_chem_s2_q09_1.png}{61.0}{0}
\qfigdef{p2_chem_s2_q09_2}{Pre-Score (2023-24) - Test 03 (22 Oct 2023)-img/p2_chem_s2_q09_2.png}{117.3}{0}
\qfigdef{p2_chem_s2_q09_3}{Pre-Score (2023-24) - Test 03 (22 Oct 2023)-img/p2_chem_s2_q09_3.png}{129.3}{0}
\qfigdef{p2_chem_s2_q09_4}{Pre-Score (2023-24) - Test 03 (22 Oct 2023)-img/p2_chem_s2_q09_4.png}{88.6}{0}
\qfigdef{p2_chem_s2_q10_1}{Pre-Score (2023-24) - Test 03 (22 Oct 2023)-img/p2_chem_s2_q10_1.png}{61.0}{0}
\qfigdef{p2_chem_s2_q10_2}{Pre-Score (2023-24) - Test 03 (22 Oct 2023)-img/p2_chem_s2_q10_2.png}{117.3}{0}
\qfigdef{p2_chem_s2_q10_3}{Pre-Score (2023-24) - Test 03 (22 Oct 2023)-img/p2_chem_s2_q10_3.png}{129.3}{0}
\qfigdef{p2_chem_s2_q10_4}{Pre-Score (2023-24) - Test 03 (22 Oct 2023)-img/p2_chem_s2_q10_4.png}{95.0}{0}
\qfigdef{p2_chem_s2_q10_5}{Pre-Score (2023-24) - Test 03 (22 Oct 2023)-img/p2_chem_s2_q10_5.png}{132.0}{0}
\qfigdef{p2_chem_s2_q10_6}{Pre-Score (2023-24) - Test 03 (22 Oct 2023)-img/p2_chem_s2_q10_6.png}{112.5}{0}
\qfigdef{p2_chem_s2_q10_7}{Pre-Score (2023-24) - Test 03 (22 Oct 2023)-img/p2_chem_s2_q10_7.png}{125.8}{0}
\qfigdef{p2_chem_s2_q10_8}{Pre-Score (2023-24) - Test 03 (22 Oct 2023)-img/p2_chem_s2_q10_8.png}{133.0}{0}
\qfigdef{p2_chem_s2_q10_9}{Pre-Score (2023-24) - Test 03 (22 Oct 2023)-img/p2_chem_s2_q10_9.png}{113.1}{0}
\qfigdef{p2_chem_s2_q11_1}{Pre-Score (2023-24) - Test 03 (22 Oct 2023)-img/p2_chem_s2_q11_1.png}{132.0}{0}
\qfigdef{p2_chem_s2_q11_2}{Pre-Score (2023-24) - Test 03 (22 Oct 2023)-img/p2_chem_s2_q11_2.png}{112.5}{0}
\qfigdef{p2_chem_s2_q11_3}{Pre-Score (2023-24) - Test 03 (22 Oct 2023)-img/p2_chem_s2_q11_3.png}{125.8}{0}
\qfigdef{p2_chem_s2_q11_4}{Pre-Score (2023-24) - Test 03 (22 Oct 2023)-img/p2_chem_s2_q11_4.png}{133.0}{0}
\qfigdef{p2_chem_s2_q11_5}{Pre-Score (2023-24) - Test 03 (22 Oct 2023)-img/p2_chem_s2_q11_5.png}{115.8}{0}
\qfigdef{p2_chem_s3_q01_1}{Pre-Score (2023-24) - Test 03 (22 Oct 2023)-img/p2_chem_s3_q01_1.png}{132.2}{0}
\qfigdef{p2_chem_s3_q01_2}{Pre-Score (2023-24) - Test 03 (22 Oct 2023)-img/p2_chem_s3_q01_2.png}{89.4}{0}
\qfigdef{p2_chem_s3_q01_3}{Pre-Score (2023-24) - Test 03 (22 Oct 2023)-img/p2_chem_s3_q01_3.png}{62.0}{0}
\qfigdef{p2_chem_s3_q01_4}{Pre-Score (2023-24) - Test 03 (22 Oct 2023)-img/p2_chem_s3_q01_4.png}{48.9}{0}
\qfigdef{p2_chem_s3_q03_1}{Pre-Score (2023-24) - Test 03 (22 Oct 2023)-img/p2_chem_s3_q03_1.png}{148.0}{0}
\qfigdef{p2_chem_s3_q03_2}{Pre-Score (2023-24) - Test 03 (22 Oct 2023)-img/p2_chem_s3_q03_2.png}{132.2}{0}
\qfigdef{p2_chem_s3_q03_3}{Pre-Score (2023-24) - Test 03 (22 Oct 2023)-img/p2_chem_s3_q03_3.png}{71.1}{0}
\qfigdef{p2_chem_s3_q04_1}{Pre-Score (2023-24) - Test 03 (22 Oct 2023)-img/p2_chem_s3_q04_1.png}{289.0}{1}
\qfigdef{p2_chem_s3_q04_2}{Pre-Score (2023-24) - Test 03 (22 Oct 2023)-img/p2_chem_s3_q04_2.png}{203.2}{0}
\qfigdef{p2_chem_s3_q05_1}{Pre-Score (2023-24) - Test 03 (22 Oct 2023)-img/p2_chem_s3_q05_1.png}{95.4}{0}
\qfigdef{p2_chem_s3_q05_2}{Pre-Score (2023-24) - Test 03 (22 Oct 2023)-img/p2_chem_s3_q05_2.png}{119.9}{0}
\qfigdef{p2_chem_s3_q05_3}{Pre-Score (2023-24) - Test 03 (22 Oct 2023)-img/p2_chem_s3_q05_3.png}{171.6}{0}
\qfigdef{p2_chem_s3_q05_4}{Pre-Score (2023-24) - Test 03 (22 Oct 2023)-img/p2_chem_s3_q05_4.png}{198.1}{0}
\qfigdef{p2_chem_s3_q05_5}{Pre-Score (2023-24) - Test 03 (22 Oct 2023)-img/p2_chem_s3_q05_5.png}{245.2}{0}
\qfigdef{p2_chem_s3_q05_6}{Pre-Score (2023-24) - Test 03 (22 Oct 2023)-img/p2_chem_s3_q05_6.png}{245.2}{0}
\qfigdef{p2_chem_s3_q06_1}{Pre-Score (2023-24) - Test 03 (22 Oct 2023)-img/p2_chem_s3_q06_1.png}{100.6}{0}
% ---- body skeleton ----
% \begin{document}
% \begin{center}{\Large\bfseries JEE (Main+Advanced) Score-II --- Test NN (DD Mon YYYY)}\end{center}
% \begin{multicols}{2}
%   \papersec{1}   (only for two-paper tests)
%   \partsec{1}{PHYSICS} \sect{Section-I (i): Single Correct} \qn{1} ... \qfig{p1_phy_q01_1} ...
% \end{multicols}
% \end{document}

\begin{document}
\begin{center}{\Large\bfseries Pre-Score (2023-24) --- Test 03 (22 Oct 2023)}\end{center}
\begin{multicols}{2}
\papersec{1}
\partsec{1}{PHYSICS}
\sect{SECTION-I(i)}
\qn{1} A black thin disc of radius 1m is kept 0.75 m away from a point source of power $56.7\,\pi$w. If the disc is highly conducting, what is its temperature in steady state? Surrounding is at 0 K.
\qfig{p1_phy_s1_q01_1}
\begin{qoptions}
\opt{A}{100 K}
\opt{B}{141 K}
\opt{C}{70.5 K}
\opt{D}{282 K}
\end{qoptions}
\qn{2} A non-linear triatomic gas is filled inside a vessel. If `$\alpha$' fraction of moles dissociate into individual atoms, then average degree of freedom for the mixture is : (neglect vibrational degrees of freedom)
\begin{qoptions}
\opt{A}{$\dfrac{3\alpha+6}{\alpha+1}$}
\opt{B}{$\dfrac{\alpha+6}{2\alpha+1}$}
\opt{C}{$\dfrac{3\alpha+6}{\alpha+2}$}
\opt{D}{$\dfrac{3\alpha+6}{2\alpha+1}$}
\end{qoptions}
\qn{3} An ideal gas goes through a cycle. Starting from state A, it is first heated at constant volume up to state B, then it expands to state C without a change in temperature and is finally isobarically compressed again to initial state A. Use p, V and T to denote the pressure, volume and temperature of the gas. Which of the following graphs represent the cycle correctly?
\qfig{p1_phy_s1_q03_1}
\begin{qoptions}
\opt{A}{Graph I and II only.}
\opt{B}{Graph I and III only.}
\opt{C}{Graph II and III only.}
\opt{D}{All three graphs.}
\end{qoptions}
\qn{4} On a cold winter's day, there are three identical, non-insulated wooden boxes in front of the house, each filled with the same amount of ice at a temperature of 0.0 $^\circ$C. To melt the ice, an electric heating element is placed in each of the boxes. The heating elements are identical, but are operated with different voltages. In the first box, the heating element is operated with a voltage of 80 V. All of the ice in the box will then melt in 20.0 minutes. A voltage of 120 V is applied to the heating element of the second box, whereupon the ice completely melts in just 4.0 minutes. In the third box, a voltage of 40 V is used for the heating element. The heating elements are designed in such a way that they heat the entire mass of ice in the respective box at the same time. Assume that the meltwater is not being heated by the heating element. Which of the following statements is then true for the melting of the ice in the third box?
\begin{qoptions}
\opt{A}{It takes about 80 minutes to melt all of the ice in the third box.}
\opt{B}{It takes about 100 minutes to melt all of the ice in the third box.}
\opt{C}{It will take about 130 minutes to melt all of the ice in the third box.}
\opt{D}{With the voltage used, it is not possible to melt all of the ice.}
\end{qoptions}
\sect{SECTION-I(ii)}
\qn{5} A process for monoatomic gas is defined by $V=aT^2$, where $a$ is a constant. In this process temperature of 1 mole of a gas is increased by 100 K; then ($V$ = Volume, $T$ = Temperature)
\begin{qoptions}
\opt{A}{Molar heat capacity of this process is negative.}
\opt{B}{The work done by the gas in the process is 250 R.}
\opt{C}{The work done by the gas in the process is 200 R.}
\opt{D}{Heat given to the gas is 350 R.}
\end{qoptions}
\qn{6} Two spherical black bodies A and B, having radii $r_A$ and $r_B$ where $r_B=2r_A$ emit radiation with peak intensities at wavelength 400 nm and 800 nm respectively. If their temperature are $T_A$ and $T_B$ respectively in Kelvin scale, their emissive powers are $E_A$ and $E_B$, and energy emitted per second are $P_A$ and $P_B$ then
\begin{qoptions}
\opt{A}{$\dfrac{T_A}{T_B}=2$}
\opt{B}{$\dfrac{P_A}{P_B}=4$}
\opt{C}{$\dfrac{E_A}{E_B}=8$}
\opt{D}{$\dfrac{E_A}{E_B}=4$}
\end{qoptions}
\qn{7} A piece of ice having temperature 0$^\circ$C is thrown into a thermally insulated container, filled with water of volume 5 dm$^3$ at a temperature of 50$^\circ$C. Match the final amount of water in the container if after the ice melts the temperature of the system is $T\,^\circ$C?
\begin{qoptions}
\opt{A}{$T=15,\ V=\dfrac{35}{19}$ dm$^3$}
\opt{B}{$T=0,\ V=\dfrac{65}{8}$ dm$^3$}
\opt{C}{$T=5,\ V=\dfrac{130}{17}$ dm$^3$}
\opt{D}{$T=10,\ V=\dfrac{100}{9}$ dm$^3$}
\end{qoptions}
\qn{8} A vessel is completely filled by a liquid of density $\rho$ as shown in figure having acceleration $g$ in horizontal direction. After getting equilibrium condition a particle of density $\dfrac{\rho}{2}$ is released from bottom of vessel at distance `$x$' as shown in figure such that it reaches to water-air interface after time $\sqrt{\dfrac{2h}{g}}$. Assume all collisions with wall are completely elastic and walls are smooth. Which of the positions satisfy the given situation :-
\qfig{p1_phy_s2_q08_1}
\begin{qoptions}
\opt{A}{$x=0$}
\opt{B}{$x=h$}
\opt{C}{$x=\dfrac{h}{9}$}
\opt{D}{$x=\dfrac{h}{4}$}
\end{qoptions}
\qn{9} The diagram below shows a cross-sectional view of a cylindrical pipe of varying diameter.
\qfig{p1_phy_s2_q09_1}
If an ideal fluid is flowing through the pipe, find the true statement(s).
\begin{qoptions}
\opt{A}{Volume flow rate is same at A \& B}
\opt{B}{The pressure is lower at point B than at point A.}
\opt{C}{The volume flow rate is greater at point A than at point B.}
\opt{D}{The flow speed is greater at point B than at point A.}
\end{qoptions}
\qn{10} In the figure shown the system is in equilibrium. In this system we have used two light wires 1 and 2. The Young's modulus and radius of the wire 1 are two times and $2^{1/4}$ times respectively that of wire 2, then.
\qfig{p1_phy_s2_q10_1}
\begin{qoptions}
\opt{A}{Elastic energy density in wire 1 is equal to that of in wire 2}
\opt{B}{Elastic energy density in wire 1 is greater than that of in wire 2}
\opt{C}{Elastic energy density in wire 1 is less than that of in wire 2}
\opt{D}{Stress in wire 1 is greater than stress in wire 2.}
\end{qoptions}
\qn{11} An ideal heat engine has two sources of finite thermal capacity $C=500$ J/K each, with initial temperature $T_H=625$ K and $T_C=289$ K. The power of the engine decreases exponentially in time according to the law $P=P_0e^{-k\tau}$, where $P_0=960$ W is the initial power, and $\tau$ is the time elapsed.
\begin{qoptions}
\opt{A}{The temperatures of the two sources, when the engine ceases to function is 457 K.}
\opt{B}{The total work done by the engine is 32 kJ.}
\opt{C}{$k=0.6$ K/sec.}
\opt{D}{Initially, the temperature $T_H$ decreases at rate of $\dfrac{25}{7}$ K/sec.}
\end{qoptions}
\qn{12} A Thermally insulated cylindrical vessel is divided into two equal parts by a non-conducting piston that can move without friction. At the initial moment in the left and right parts of the vessel contains one mole of helium at the same temperature. Heat was supplied to the left side of the vessel using a heater. At the same time, the temperature of the helium in it increased by a small amount $\Delta T$. The temperature changes by $\Delta T_2$ on the right side of the vessel and the amount of heat $Q$ is transferred by the heater to the system.
\qfig{p1_phy_s2_q12_1}
\begin{qoptions}
\opt{A}{$\Delta T_2=\dfrac{\Delta T}{4}$}
\opt{B}{$\Delta T_2=\dfrac{\Delta T}{2}$}
\opt{C}{$Q=\dfrac{9R\Delta T}{4}$}
\opt{D}{$Q=\dfrac{15R\Delta T}{8}$}
\end{qoptions}
\sect{SECTION-II}
\qn{1} A glass bulb with volumetric expansion coefficient $\alpha=1.50\times10^{-5}\,{}^\circ\mathrm{C}^{-1}$ is weighed in liquid at temperatures $t=10^\circ$C and $t_1=20^\circ$C. The weights of the displaced liquids are $p=100.01$ g and $p_1=100$ g respectively. Find the volumetric expansion coefficient $\alpha_1$ (in the units of $10^{-5}\,{}^\circ\mathrm{C}^{-1}$) of the liquid in the temperature interval from $t$ to $t_1$.
\qn{2} A uniform tube PQR is in the form of right angle as shown in the figure. It is closed at R and has a small cross-sectional area A. It is filled with water till point P and is accelerating horizontally with acceleration ``a'' as shown. Find the acceleration $a$ (in m/s$^2$) so that pressure at R is atmospheric. (Take $g=10$ m/s$^2$)
\qfig{p1_phy_s3_q02_1}
\qn{3} A drop of mercury is present between two glass plates each of mass 50 gm. Separation between the plates is negligible as compared to spread of the mercury between the plates. What should be the mass of a block (in kg) which should be placed over the plate A so that the separation between the plates becomes half of the initial separation. Assume contact angle between mercury and glass is 180$^\circ$.
\qfig{p1_phy_s3_q03_1}
\qn{4} A rocket-powered sled is used in the training of astronauts as shown in figure. For braking, scoops are lowered to deflect water from a stationary tank of water. To what height $h$ (in m) does a sled traveling at 30 m/s deflect water? Take $\sin 14.5^\circ=0.25$
\qfig{p1_phy_s3_q04_1}
\qfig{p1_phy_s3_q04_2}
\qn{5} There's a one mole diatomic ideal gas following a thermodynamic cycle. The process consists of an isothermal compression followed by an isobaric heating, which is followed by a polytropic process and isochoric. The molar capacity of gas during polytropic process is $C=0.5R$. Find the efficiency of this cycle in percentage. Take $\ell n2=0.7$
\qfig{p1_phy_s3_q05_1}
\qn{6} A coin is frozen into a piece of ice weighing 130 g and density 900 kg/m$^3$. The coin weighs 10 g and density 9000 kg/m$^3$. This piece of ice with the coin has temperature 0 $^\circ$C, placed in a vessel containing 400 ml of water with some initial temperature $t$. Heat exchange with the surrounding can be neglected. The ice and coin initially float without touching the bottom of the vessel. Specific heat capacity of water 4200 J/(kg $\cdot\,^\circ$C), specific heat of melting of ice is 330 kJ/kg, water density 1000 kg/m$^3$, acceleration free fall 10 m/s$^2$. What should be the minimum initial water temperature so that remaining piece of ice along with a coin sank to the bottom after it attains thermal equilibrium? Express your answer in degrees celsius.
\partsec{2}{CHEMISTRY}
\sect{SECTION-I(i)}
\qn{1} Major product ``Z'' obtained in given reaction is ?
\qfig{p1_chem_s1_q01_1}
\qfig{p1_chem_s1_q01_2}
\begin{qoptions}
\opt{A}{\qfigi{p1_chem_s1_q01_A}}
\opt{B}{\qfigi{p1_chem_s1_q01_B}}
\opt{C}{\qfigi{p1_chem_s1_q01_C}}
\opt{D}{\qfigi{p1_chem_s1_q01_D}}
\end{qoptions}
\qn{2} \qfig{p1_chem_s1_q02_1}
Choose the incorrect option ?
\begin{qoptions}
\opt{A}{(S) is 4-bromoaniline}
\opt{B}{(Q) is Acetanilide}
\opt{C}{Basic strength order is : (Q) $>$ (P) $>$ (R)}
\opt{D}{Basic strength order is : (P) $>$ (R) $>$ (Q)}
\end{qoptions}
\qn{3} Which of the following reaction sequence would be the best for synthesizing the compound, 1-bromo-3-propylbenzene ?
\qfig{p1_chem_s1_q03_1}
\begin{qoptions}
\opt{A}{\qfigi{p1_chem_s1_q03_A}}
\opt{B}{\qfigi{p1_chem_s1_q03_B}}
\opt{C}{\qfigi{p1_chem_s1_q03_C}}
\opt{D}{\qfigi{p1_chem_s1_q03_D}}
\end{qoptions}
\qfullbegin
\qn{4} Product of which of the following reaction gives carbyl amine test.
\qfig{p1_chem_s1_q04_1}
\qfullsplit
\qfig{p1_chem_s1_q04_2}
\begin{qoptions}
\opt{A}{I and II}
\opt{B}{I, II, III}
\opt{C}{I and III}
\opt{D}{III only}
\end{qoptions}
\qfullend
\sect{SECTION-I(ii)}
\qn{5} Identify reactions incorrectly matched with their final products?
\begin{qoptions}
\opt{A}{\qfigi{p1_chem_s2_q05_A}}
\opt{B}{\qfigi{p1_chem_s2_q05_B}}
\opt{C}{\qfigi{p1_chem_s2_q05_C}}
\opt{D}{\qfigi{p1_chem_s2_q05_D}}
\end{qoptions}
\qn{6} \qfig{p1_chem_s2_q06_1}
Products which is / are formed in above reaction
\newsavebox{\bxB}\sbox{\bxB}{\qfigi{p1_chem_s2_q06_B}}
\begin{qoptions}
\opt{A}{\qfigi{p1_chem_s2_q06_A}}
\opt{B}{\usebox{\bxB}}
\opt{C}{\qfigi{p1_chem_s2_q06_C}}
\opt{D}{\qfigi{p1_chem_s2_q06_D}}
\end{qoptions}
\qn{7} Identify reactions correctly matched with their major products :
\begin{qoptions}
\opt{A}{\qfigi{p1_chem_s2_q07_A}}
\opt{B}{\qfigi{p1_chem_s2_q07_B}}
\opt{C}{\qfigi{p1_chem_s2_q07_C}\quad\qfigi{p1_chem_s2_q07_4}}
\opt{D}{\qfigi{p1_chem_s2_q07_D}}
\end{qoptions}
\qn{8} Which are Correct statement(s) in the following?
\begin{qoptions}
\opt{A}{Rate of sulphonation of C$_6$H$_6$ and C$_6$D$_6$ are equal.}
\opt{B}{Rate of halogenation of\par\smallskip\qfigi{p1_chem_s2_q08_1}}
\opt{C}{\qfigi{p1_chem_s2_q08_2}\ on reaction with Br$_2$ / H$_2$O forms tribromoderivative.}
\opt{D}{\qfigi{p1_chem_s2_q08_3}\ on reaction with Cl$_2$ at 500$^\circ$C will give electrophilic aromatic substitution reaction.}
\end{qoptions}
\qn{9} Which of the following will react at rate faster than benzene during aromatic electrophilic substitution?
\begin{qoptions}
\opt{A}{\qfigi{p1_chem_s2_q09_A}}
\opt{B}{\qfigi{p1_chem_s2_q09_B}}
\opt{C}{\qfigi{p1_chem_s2_q09_C}}
\opt{D}{\qfigi{p1_chem_s2_q09_D}}
\end{qoptions}
\qn{10} Select reaction(s) with correct major product :
\begin{qoptions}
\opt{A}{\qfigi{p1_chem_s2_q10_1}\ \qfigi{p1_chem_s2_q10_A}}
\opt{B}{PhCH$_2$COOH $\xrightarrow[\text{CaO}/\Delta]{\text{NaOH}}$ PhCH$_3$}
\opt{C}{\qfigi{p1_chem_s2_q10_C}}
\opt{D}{\qfigi{p1_chem_s2_q10_D}\ \qfigi{p1_chem_s2_q10_5}}
\end{qoptions}
\qn{11} \qfig{p1_chem_s2_q11_1}
Identify correct statement for above reaction sequence ?
\begin{qoptions}
\opt{A}{`Z' product is sweet smelling substance}
\opt{B}{Reaction of\par\smallskip\qfigi{p1_chem_s2_q11_2}}
\opt{C}{`Y' product librate CO$_2$ gas with aqueous NaHCO$_3$ and also soluble in aqueous NaOH}
\opt{D}{`Z' product is resolvable}
\end{qoptions}
\qn{12} Which of the following conditions or reaction sequence will NOT give acetophenone as the major product ?
\begin{qoptions}
\opt{A}{\qfigi{p1_chem_s2_q12_1}}
\opt{B}{\qfigi{p1_chem_s2_q12_2}}
\opt{C}{\qfigi{p1_chem_s2_q12_C}}
\opt{D}{\qfigi{p1_chem_s2_q12_D}}
\end{qoptions}
\sect{SECTION-II}
\qn{1} How many compounds will give Anesthesia product with Cl$_2$ / NaOH :
\qfig{p1_chem_s3_q01_1}
\qfig{p1_chem_s3_q01_2}
\qfig{p1_chem_s3_q01_3}
\qfullbegin
\qn{2} Find total number of reaction in which acetic acid is one of product.
\qfig{p1_chem_s3_q02_1}
\qfullsplit
\qfig{p1_chem_s3_q02_3}
\qfig{p1_chem_s3_q02_4}
\qfig{p1_chem_s3_q02_5}
\qn{3} Total number of reactions which give aromatic product as a major final product?
\qfig{p1_chem_s3_q03_1}
\qfullsplit
\qfig{p1_chem_s3_q03_2}
\qfig{p1_chem_s3_q03_3}
\qfig{p1_chem_s3_q03_4}
\qfullend
\qn{4} How many reactions will give 1$^\circ$ amine as final product :
\qfig{p1_chem_s3_q04_1}
\qfig{p1_chem_s3_q04_2}
\qfig{p1_chem_s3_q04_3}
\par\noindent(vii) CH$_3-$CH$_2-$C$\equiv$N $\xrightarrow{\text{Na/EtOH}}$
\qn{5} How many of the following can be used to differentiate between CH$_3-$CH$_2-$OH and \qfigi{p1_chem_s3_q05_1}

(a) NaOI \quad (b) Neutral FeCl$_3$ \quad (c) Azo-dye test \quad (d) Libermann nitroso test \quad (e) Na \quad (f) NaHCO$_3$
\qn{6} Find number of compound which are soluble in NaOH as well as NaHCO$_3$?
\qfig{p1_chem_s3_q06_1}
\qfig{p1_chem_s3_q06_2}
\par\noindent(v) CH$_3$OH \quad (vi) \qfigi{p1_chem_s3_q06_3}
\qfig{p1_chem_s3_q06_4}
\partsec{3}{MATHEMATICS}
\sect{SECTION-I(i)}
\qn{1} $\cos\dfrac{\pi}{7}-\cos\dfrac{2\pi}{7}+\cos\dfrac{3\pi}{7}$ is
\begin{qoptions}
\opt{A}{$\dfrac12$}
\opt{B}{$-\dfrac32$}
\opt{C}{$-\dfrac14$}
\opt{D}{$\dfrac18$}
\end{qoptions}
\qn{2} Number of solution of $\dfrac{\tan\theta}{\sqrt3}+\dfrac{\tan2\theta}{\sqrt3}+\tan\theta\tan2\theta=1$ in $\theta\in(0,2\pi)$
\begin{qoptions}
\opt{A}{0}
\opt{B}{1}
\opt{C}{4}
\opt{D}{6}
\end{qoptions}
\qn{3} If lines $L_1,L_2$ are the member of family $(x-2)+\lambda(y+1)=0$ and $(x-5)+\mu(y-3)=0$ respectively (where $\lambda$ and $\mu$ are parameter). Then maximum possible perpendicular distance between $L_1$ and $L_2$ is
\begin{qoptions}
\opt{A}{$\dfrac23$}
\opt{B}{3}
\opt{C}{$\dfrac72$}
\opt{D}{5}
\end{qoptions}
\qn{4} P is a point on the line $y+2x=1$ and Q and R are two points on the line $3y+6x=6$ such that the triangle PQR is an equilateral triangle. Area of this triangle is equal to
\begin{qoptions}
\opt{A}{$\dfrac{1}{15}$}
\opt{B}{$\dfrac{1}{5\sqrt3}$}
\opt{C}{$\dfrac{1}{3\sqrt5}$}
\opt{D}{$\dfrac{1}{2\sqrt{15}}$}
\end{qoptions}
\sect{SECTION-I(ii)}
\qn{5} If $\sin x\cos^5x>\cos x\sin^5x$ then $x$ lies in
\begin{qoptions}
\opt{A}{$\Bigl(0,\dfrac{\pi}{8}\Bigr)$}
\opt{B}{$\Bigl(\dfrac{\pi}{8},\dfrac{\pi}{4}\Bigr)$}
\opt{C}{$\Bigl(\dfrac{\pi}{2},\dfrac{5\pi}{8}\Bigr)$}
\opt{D}{$\Bigl(\dfrac{5\pi}{8},\dfrac{3\pi}{4}\Bigr)$}
\end{qoptions}
\qn{6} If which of the following is/are correct
\begin{qoptions}
\opt{A}{If $\cos\theta=\dfrac{a+\frac1a}{2}$ then $\cos(3\theta)=\dfrac{a^3+\frac1{a^3}}{2}$}
\opt{B}{If $\cos x+\cos y+\cos z=0$ then $\cos3x+\cos3y+\cos3z=12\cos x\cos y\cos z$}
\opt{C}{$\dfrac{\cos9^\circ+\sin9^\circ}{\cos9^\circ-\sin9^\circ}=\tan(54^\circ)$}
\opt{D}{If $\sin\theta$ is geometric mean of $\sin\phi$ and $\cos\phi$ then $\cos2\theta=2\cos^2\Bigl(\dfrac{\pi}{4}+\phi\Bigr)$}
\end{qoptions}
\qn{7} Let $y=\dfrac{1}{\Bigl(\sin^4x+\cos^4x+\frac{\sin2x}{2}+1\Bigr)}$ then $y$ can be
\begin{qoptions}
\opt{A}{$\dfrac14$}
\opt{B}{$\dfrac{8}{17}$}
\opt{C}{1}
\opt{D}{0.8}
\end{qoptions}
\qn{8} All values of $x\in\Bigl(0,\dfrac{\pi}{2}\Bigr)$ such that $\dfrac{\sqrt3-1}{\sin x}+\dfrac{\sqrt3+1}{\cos x}=4\sqrt2$ are
\begin{qoptions}
\opt{A}{$\dfrac{\pi}{15}$}
\opt{B}{$\dfrac{\pi}{12}$}
\opt{C}{$\dfrac{11\pi}{36}$}
\opt{D}{$\dfrac{3\pi}{10}$}
\end{qoptions}
\qn{9} If $a\cos x+\cot x+1=\operatorname{cosec}x$ then the possible values of $x$ can be
\begin{qoptions}
\opt{A}{$x=n\pi$, $a\in R$, $n\in I$}
\opt{B}{$x=(4n+1)\dfrac{\pi}{2}$, $a\in R$, $n\in I$}
\opt{C}{$x=\dfrac12\sin^{-1}\Bigl(\dfrac{4(a+1)}{a^2}\Bigr)$, $a\in\bigl(-\infty,2-2\sqrt2\bigr]\cup\bigl[2+2\sqrt2,\infty\bigr)-\{-1\}$}
\opt{D}{none of these}
\end{qoptions}
\qn{10} The equation of the circle which touches the axes of coordinates and the line $\dfrac x3+\dfrac y4=1$ and whose centre lies in the first quadrant is $x^2+y^2-2cx-2cy+c^2=0$ where $c$ is
\begin{qoptions}
\opt{A}{1}
\opt{B}{2}
\opt{C}{3}
\opt{D}{6}
\end{qoptions}
\qn{11} If $\tan\theta_1,\tan\theta_2,\tan\theta_3$ are the real roots of the equation $x^3-(a+1)x^2+(b-a)x-b=0$, where $\theta_1+\theta_2+\theta_3\in(0,\pi)$, then $\theta_1+\theta_2+\theta_3$ is equal to
\begin{qoptions}
\opt{A}{$\dfrac{\pi}{2}$}
\opt{B}{$\dfrac{\pi}{4}$}
\opt{C}{$\dfrac{3\pi}{4}$}
\opt{D}{$\pi$}
\end{qoptions}
\qn{12} Two circles with radii $\sqrt2$ cm and 1 cm meet at a point A, whose centres lie on $x$-axis. The distance between their centres is 2 cm. The chord AC of the larger circle cuts the smaller circle at a point B and is bisected by that point. The length of chord AC is not equal to
\begin{qoptions}
\opt{A}{$\dfrac{\sqrt5}{2}$ cm}
\opt{B}{$\sqrt{\dfrac72}$ cm}
\opt{C}{$\dfrac{\sqrt3}{2}$ cm}
\opt{D}{$\sqrt{\dfrac52}$ cm}
\end{qoptions}
\sect{SECTION-II}
\qn{1} Let P is lying on $x$-axis and B$(-2,4)$, C$(-3,5)$ are given then the number of possible points P such that $PA+PB=7$ is/are
\qn{2} In a triangle ABC, the least values of $\prod\operatorname{cosec}\Bigl(\dfrac A2\Bigr)$ and $\sum\operatorname{cosec}^2A$ are $p$ and $q$ respectively, then $(p-q)$ is equal to
\qn{3} A diagonal of rhombus ABCD is a member of both the family of lines $(x+y-1)+\lambda_1(2x+3y-2)=0$ and $(x-y+2)+\lambda_2(2x-3y+5)=0$ where $\lambda_1,\lambda_2\in R$ and one of the vertex of the rhombus is $(3,2)$. If area of the rhombus is $12\sqrt5$ square units then find the length of semi longer diagonal of the rhombus
\qn{4} Number of solution of the equation $\sin^4x-\cos^2x\sin x+2\sin^2x+\sin x=0$ in $0\le x\le3\pi$ is
\qn{5} The centres of the three circles A, B and C are collinear with the centre of circle B lying between the centres of circle A and C. Circles A and C both touch externally to circle B and the three circles share a common tangent line. Given that circle A has radius 12 and circle B has radius 42 the radius of circle C, is equal to
\qn{6} Curves $ax^2+2hxy+by^2-2gx-2fy+c=0$ and $a'x^2-2hxy+(a'+a-b)y^2-2g'x-2f'y+c=0$ intersect at four concyclic points A,B, C are D. If P is the point $\Bigl(\dfrac{g'+g}{a'+a},\dfrac{f'+f}{a'+a}\Bigr)$, the value of $\dfrac{PA^2+PB^2+PC^2}{PD^2}$ is
\papersec{2}
\partsec{1}{PHYSICS}
\sect{SECTION-I(i)}
\qn{1} The tips of two identical syringes were connected using a short rubber tube. Initially, the volume of air in each of the syringes was $V=8$ ml, the pressure inside system was equal to atmospheric pressure = 100 kPa. Cross-sectional area syringe is $S=3$ cm$^2$. Holding the syringe bodies motionless, the piston of syringe No. 1 is pressed so that it begins to move very slowly with a constant speed. On the piston of syringe No. 2, no additional force is applied. At the point in time when the volume of air in syringe No. 1 was $V_1=6$ ml, the volume of air in syringe No. 2 was $V_2=9$ ml. There is no stop anywhere in the syringes.
\qfig{p2_phy_s1_q01_1}
\begin{qoptions}
\opt{A}{The piston of syringe 2 remained motionless for some time and then started moving}
\opt{B}{The piston of syringe 2 moved initially and then abruptly stopped}
\opt{C}{The pressure in the syringes in the final position is approximately 106.6 kPa}
\opt{D}{The pressure in the syringes in the final position is approximately 110.2 kPa}
\end{qoptions}
\qn{2} The ends of a metallic bar are maintained at different temperature. In the steady state the temperature variation along the length of the bar is as shown in the figure. This could possibly be due to :
\qfig{p2_phy_s1_q02_1}
\begin{qoptions}
\opt{A}{No heat loss to the surroundings along the length of the bar. The area of cross-section is decreasing as we move from hot end to cold end.}
\opt{B}{No heat loss to the surroundings along the length of the bar. The coefficient of thermal conductivity is decreasing as we moved from hot end to cold end.}
\opt{C}{No heat loss to the surroundings along the length of the bar. the area of cross-section is increasing as we move from hot end to cold end.}
\opt{D}{The bar is uniform but heat is being lost from the sides along the length of the bar to the surroundings.}
\end{qoptions}
\qn{3} A cubical piece of metal floats on mercury (Density of metal is half the density of mercury) such that its flat face is horizontal. The coefficient of volume expansion of metal and mercury are $\gamma$ and $2\gamma$ respectively. If the temperature of both mercury and metal are increased by an amount $\Delta T$,
\begin{qoptions}
\opt{A}{The fraction of volume inside the mercury increases by $2\gamma\,\Delta T$.}
\opt{B}{The Buoyant force remains the same.}
\opt{C}{The fraction of height of the cube above mercury increases by $\gamma\,\Delta T$.}
\opt{D}{The area of the cube increases by a factor of $\Bigl(1+\dfrac{2\gamma\Delta T}{3}\Bigr)$}
\end{qoptions}
\qn{4} Three ice cubes are floating in each glass containing water. The ice cube in glass 1 has an air bubble in it. The ice cube in glass 2 has a cavity inside filled with water and the third ice cube floats in glass 3 with an aluminium core. What can be said about the water level in the glasses after ice has melted ?
\qfig{p2_phy_s1_q04_1}
\begin{qoptions}
\opt{A}{The water in glass 1 has risen.}
\opt{B}{The water level in glass 3 has got down.}
\opt{C}{The water level in glass 2 is unchanged.}
\opt{D}{The water level in all glass is unchanged.}
\end{qoptions}
\qn{5} A soap bubble is very slowly blown at the end of a glass tube by a pump which supplies a fixed volume of air every minute. Here we draw the graphs of excess pressure $p$ with time $t$, surface energy $E$ with time $t$, choose the correct option(s)
\begin{qoptions}
\opt{A}{\qfigi{p2_phy_s1_q05_A}}
\opt{B}{\qfigi{p2_phy_s1_q05_B}}
\opt{C}{\qfigi{p2_phy_s1_q05_C}}
\opt{D}{\qfigi{p2_phy_s1_q05_D}}
\end{qoptions}
\qn{6} The velocity profile of a liquid flowing through the pipe is approximated by the truncated conical distribution as shown. The liquid has a density of $\rho=900$ kg/m$^3$. The coefficient of viscosity is 0.01 Ns/m. Assume Newton's law.
\qfig{p2_phy_s1_q06_1}
\begin{qoptions}
\opt{A}{The shear stress at the walls is same as the shear stress at a distance of 150 mm from the axis.}
\opt{B}{The shear stress at the axis is zero.}
\opt{C}{The force of fluid friction acting on the walls per unit length is $\dfrac{12\pi}{100}$ N/m.}
\opt{D}{The rate of volume flow through the pipe is $\dfrac{7\pi}{100}$ m$^3$/s.}
\end{qoptions}
\qn{7} A partly full water tank admits water at 20$^\circ$C and 85 N/s weight flow while ejecting water on the other side at 5500 cm$^3$/s. The air picket in the tank has a vent at the top and is at 20$^\circ$C and 1 atm. If the fluids are approximately in-compressible, how much air is passing through the vent ? In which direction ? Density of water = 1000 kg/m$^3$, $g=10$ m/s$^2$.
\begin{qoptions}
\opt{A}{Air is flowing in}
\opt{B}{Air is flowing out.}
\opt{C}{$\Bigl|\dfrac{dV}{dt}\Bigr|=3000$ cm$^3$/s}
\opt{D}{$\Bigl|\dfrac{dV}{dt}\Bigr|=4650$ cm$^3$/s}
\end{qoptions}
\qn{8} A homogeneous block with a mass $m=100$ kg hangs on three light vertical wires of equal natural lengths arranged symmetrically (Figure). Find the tension in the wires if the middle wire is of steel and the other two are of copper. All the wires have the same cross section. Consider the modulus of elasticity of steel to be double that of copper. ($g=10$ m/s$^2$)
\qfig{p2_phy_s1_q08_1}
\begin{qoptions}
\opt{A}{Tension in steel wire = 250 N}
\opt{B}{Tension in copper wire = 500 N}
\opt{C}{Tension in steel wire = 500 N}
\opt{D}{Tension in copper wire = 250 N}
\end{qoptions}
\sect{SECTION-I(ii)}
\qn{9} A solid right angled prism of dimensions shown here is submerged into a liquid such that its base is horizontal. The density of the liquid is 1gm/cc and the density of the prism is $\rho$. List$-$I gives the details of the face of the prism and List$-$II gives the force exerted by the liquid on it. Assume that the top vertex of the prism is at the surface of liquid. Neglect atmospheric pressure.
\qfig{p2_phy_s2_q09_1}
\begin{center}\small\renewcommand{\arraystretch}{1.3}\setlength{\tabcolsep}{3pt}
\begin{tabular}{|l@{\hspace{3pt}}p{2.7cm}|l@{\hspace{3pt}}l|}
\hline \multicolumn{2}{|c|}{List$-$I} & \multicolumn{2}{c|}{List$-$II} \\\hline
(I) & Inclined face & (P) & 10 N \\\hline
(II) & vertical triangular face & (Q) & 5 N \\\hline
(III) & Base & (R) & 10/3 N \\\hline
(IV) & Flat vertical rectangular face & (S) & 5/3 N \\\hline
\multicolumn{2}{|c|}{} & (T) & 20/3 N \\\hline
\multicolumn{2}{|c|}{} & (U) & $5\sqrt2$ N \\\hline
\end{tabular}
\end{center}
\begin{qoptions}
\opt{A}{I $\to$ Q; II $\to$ R; III $\to$ P; IV $\to$ S}
\opt{B}{I $\to$ U; II $\to$ S; III $\to$ P; IV $\to$ Q}
\opt{C}{I $\to$ U; II $\to$ R; III $\to$ P; IV $\to$ Q}
\opt{D}{I $\to$ U; II $\to$ T; III $\to$ P; IV $\to$ S}
\end{qoptions}
\qn{10} A solid right angled prism of dimensions shown here is submerged into a liquid such that its base is horizontal. The density of the liquid is $\rho$ and the density of the prism is 1 gm/cc. List$-$I gives the details of the density $\rho$ and List$-$II gives the acceleration of COM of the prism and angular acceleration of prism. Assume that the top vertex of the prism is at the surface of liquid. Neglect atmospheric pressure.
\qfig{p2_phy_s2_q10_1}
\begin{center}\small\renewcommand{\arraystretch}{1.3}\setlength{\tabcolsep}{3pt}
\begin{tabular}{|l@{\hspace{3pt}}p{2.7cm}|l@{\hspace{3pt}}p{2.5cm}|}
\hline \multicolumn{2}{|c|}{List$-$I} & \multicolumn{2}{c|}{List$-$II} \\\hline
(I) & 1 gm/cc & (P) & Acceleration is zero \\\hline
(II) & 2 gm/cc & (Q) & Angular acceleration is zero \\\hline
(III) & 0.5 gm/cc from top till depth of 10 cm and 2 gm/cc below the prism & (R) & Acceleration is upwards \\\hline
(IV) & 1gm/cc from top till depth of 10 cm and 2 gm/cc below the prism & (S) & Angular acceleration is clockwise \\\hline
\multicolumn{2}{|c|}{} & (T) & Angular acceleration is anticlockwise \\\hline
\multicolumn{2}{|c|}{} & (U) & Acceleration is downwards \\\hline
\end{tabular}
\end{center}
\begin{qoptions}
\opt{A}{I $\to$ P,Q; II $\to$ Q,R; III $\to$ Q,U; IV $\to$ P,Q}
\opt{B}{I $\to$ P,Q; II $\to$ Q,U; III $\to$ Q,R; IV $\to$ P,Q}
\opt{C}{I $\to$ P,Q; II $\to$ Q,R; III $\to$ S,U; IV $\to$ T,Q}
\opt{D}{I $\to$ P,Q; II $\to$ Q,R; III $\to$ T,U; IV $\to$ T,U}
\end{qoptions}
\qn{11} In the archives of Lord Kelvin, researchers discover a graph which is shown below. It refers to processes carried out on a monoatomic ideal gas in the direction shown. But it is not known that what is plotted on $x$ axis and what on $y$ axis. There is a detailed discussion amongst researchers as to what is the pressure change, volume change and temperature change in each of the processes. List$-$I gives the suggested variables for $x$ and $y$ axis. List$-$II gives the effect on the parameters shown. Match the appropriate entries.
\qfig{p2_phy_s2_q11_1}
\begin{center}\small\renewcommand{\arraystretch}{1.3}\setlength{\tabcolsep}{3pt}
\begin{tabular}{|l@{\hspace{3pt}}p{2.7cm}|l@{\hspace{3pt}}p{2.3cm}|}
\hline \multicolumn{2}{|c|}{List$-$I} & \multicolumn{2}{c|}{List$-$II} \\\hline
(I) & Temperature on $x$ axis and pressure on $y$ axis & (P) & Pressure is increasing \\\hline
(II) & Temperature on $x$ axis and volume on $y$ axis & (Q) & Pressure is decreasing \\\hline
(III) & Volume on $x$ axis and pressure on $y$ axis & (R) & Volume is increasing \\\hline
(IV) & entropy on $x$ axis and temperature on $y$ axis & (S) & Volume is decreasing \\\hline
\multicolumn{2}{|c|}{} & (T) & Temperature is increasing \\\hline
\multicolumn{2}{|c|}{} & (U) & Temperature is decreasing \\\hline
\end{tabular}
\end{center}
\begin{qoptions}
\opt{A}{I $\to$ P,T; II $\to$ R,T; III $\to$ P,R,T; IV $\to$ Q,R,T}
\opt{B}{I $\to$ P,T; II $\to$ R,T; III $\to$ P,R,T; IV $\to$ T}
\opt{C}{I $\to$ P,T; II $\to$ Q,T; III $\to$ P,R,T; IV $\to$ Q,S,T}
\opt{D}{I $\to$ P,U; II $\to$ R,U; III $\to$ P,R,U; IV $\to$ Q,S,T}
\end{qoptions}
\qn{12} In the archives of Lord Kelvin, researchers discover a graph which is shown below. It refers to process carried out on a monoatomic ideal gas in the direction shown. But it is not known that what is plotted on $x$ axis and what on $y$ axis. There is a detailed discussion amongst researchers as to what is the heat given, work done by the gas and internal energy change in each of the processes. List$-$I gives the suggested variables for $x$ and $y$ axis. List$-$II gives the effect on the parameters shown. Match the appropriate entries.
\qfig{p2_phy_s2_q12_1}
\begin{center}\small\renewcommand{\arraystretch}{1.3}\setlength{\tabcolsep}{3pt}
\begin{tabular}{|l@{\hspace{3pt}}p{2.7cm}|l@{\hspace{3pt}}p{2.3cm}|}
\hline \multicolumn{2}{|c|}{List$-$I} & \multicolumn{2}{c|}{List$-$II} \\\hline
(I) & Temperature on $x$ axis and pressure on $y$ axis & (P) & Work is done by the gas \\\hline
(II) & Temperature on $x$ axis and volume on $y$ axis & (Q) & Work is done on the gas \\\hline
(III) & Volume on $x$ axis and pressure on $y$ axis & (R) & Heat is given to the gas \\\hline
(IV) & entropy on $x$ axis and temperature on $y$ axis & (S) & Heat is extracted from the gas \\\hline
\multicolumn{2}{|c|}{} & (T) & Internal energy of the gas increases \\\hline
\multicolumn{2}{|c|}{} & (U) & Internal energy of the gas decreases \\\hline
\end{tabular}
\end{center}
\begin{qoptions}
\opt{A}{I $\to$ Q,S,U; II $\to$ Q,S,T; III $\to$ P,R,U; IV $\to$ T}
\opt{B}{I $\to$ R,T; II $\to$ Q,R,T; III $\to$ Q,R,T; IV $\to$ Q,R,T}
\opt{C}{I $\to$ R,T; II $\to$ P,R,T; III $\to$ P,R,T; IV $\to$ P,R,T}
\opt{D}{I $\to$ R,T; II $\to$ P,R,T; III $\to$ P,R,T; IV $\to$ R,T}
\end{qoptions}
\sect{SECTION-II}
\qn{1} The heat machine of the room which works as a split-system operates on the inverse Carnot cycle, and it can be assumed that the room and outdoor temperatures correspond to the temperatures on the isotherms of the cycle. In summer, when the temperature outside the window is $+27^\circ$ C. The split system, working in air conditioning mode (as a refrigeration unit). keeps the room temperature at $+17\,^\circ$C and consumes from the power grid an average power $N_1$. At the beginning of winter, when the temperature outside drops to $-3^\circ$C. the split system switches to heat pump mode and maintains the same room temperature of $+17^\circ$ C as in the summer, by consuming average power $N_2$. We can assume that the heat flow (through the windows and walls) is proportional to the temperature difference between the room and outside, with same coefficient of proportionality in summer and winter. Find the ratio of powers $n=N_2/N_1$.
\qn{2} Consider a long, solid cylinder as shown in the figure, which could represent a current-carrying wire or a fuel element in a nuclear reactor. For steady state conditions, the rate at which heat is generated within the cylinder must equal the rate at which heat is convected from the surface of the cylinder to a moving fluid. This condition allows the surface temperature to be maintained at a fixed value of $T_S$. The rate of heat generated per unit volume is 20 kW/ m$^3$, the coefficient of thermal conductivity is 5 W/mK, what is the temperature (in K) at the distance $r=R_0/2$ from the axis. Take $R_0=20$ cm. The temperature $T_s$ at the surface is 300 K.
\qfig{p2_phy_s3_q02_1}
\qn{3} A liquid takes 5 minutes to cool from 80$^\circ$C to 50$^\circ$C How much time (in mins.) will it take to cool from 60$^\circ$C to 30$^\circ$C? The temperature of surrounding is 20$^\circ$C. Assume newton's law. (Take $\ell n\,2=0.69$, $\ell n3=1.10$)
\qn{4} It is known that in a pan kept on the stove, 1 liter of water is heated by 10 degrees in 42 s. The burner power is constant, heat capacity of pan in negligible, heat loss is negligible. Specific heat capacity of water 4200 J/(kg $^\circ$C), specific heat capacity of ice 2100 J/(kg $^\circ$C), specific heat of melting of ice 330 kJ/kg. How long does it take to heat on the same burner to boiling point of water a mixture of 1 liter of water and 3 kilograms of ice? Initially, the mixture is at the melting temperature of ice. Give your answer in minutes.
\qfullbegin
\qn{5} A long plate of width $\ell$ is brought into contact with the liquid surface. Then the plate is pulled up. Find the value of the force exerted by us per unit length (in N/m) of the plate. Liquid density $\rho$ and surface tension ($\sigma$). $\lambda$ = mass per unit length.\\
(Take $\lambda=10$ gm/m, $\rho=10^3$ kg/m$^3$, $\sigma=0.5$ N/m, $x=2$mm, $\ell=10$ cm, $g=10$ m/s$^2$.)
\qfig{p2_phy_s3_q05_1}
\qfullend
\qn{6} A sprinkler has four outlets as shown. All of them have uniform cross sectional area of 10 cm$^2$. Water of density 1000 kg/m$^3$ enters the system at the centre as shown at a flow rate of 200 cc/sec. The ejected water exerts a torque so that in the steady state, the sprinkler rotates with a constant angular velocity $\omega$ as shown. There in no external torque. Neglect friction. Find the steady state angular velocity (in rad/s) of the sprinkler. $r=5$ cm, $b=1$ mm.
\qfig{p2_phy_s3_q06_1}
\partsec{2}{CHEMISTRY}
\sect{SECTION-I(i)}
\qn{1} For the given conversion correct option among the following is
\qfig{p2_chem_s1_q01_1}
\qfig{p2_chem_s1_q01_2}
\begin{qoptions}
\opt{A}{Maximum number of RMgX molecule that can be consumed by one molecule of A is 6}
\opt{B}{Total number of CO$_2$ molecule release on heating one molecule of C is 2}
\opt{C}{Final product D is a ketone}
\opt{D}{Final product D is an aldehyde and can consume only 1 molecule of RMgX}
\end{qoptions}
\qn{2} Which of the following has benzene ring where both the groups are deactivating in the nature?
\begin{qoptions}
\opt{A}{\qfigi{p2_chem_s1_q02_A}}
\opt{B}{\qfigi{p2_chem_s1_q02_B}}
\opt{C}{\qfigi{p2_chem_s1_q02_C}}
\opt{D}{\qfigi{p2_chem_s1_q02_D}}
\end{qoptions}
\qn{3} Among the following, which statement is not correct ?
\begin{qoptions}
\opt{A}{When phenol is reacted with CHCl$_3$ + NaOH then product is salicylaldehyde.}
\opt{B}{Ester is formed as major product when acetic acid treated with P$_2$O$_5$/$\Delta$}
\opt{C}{Fehling's solution is a good reagent to detect aromatic aldehydes}
\opt{D}{Both aldehyde and ketone can react with 2,4-dinitrophenylhydrazine reagent}
\end{qoptions}
\qn{4} CH$_3$(CH$_2$)$_{14}$COOH $\xrightarrow{\text{LiAlH}_4}$ W $\xrightarrow{\text{HCl}}$ X\allowbreak\ $\xrightarrow[\text{2.Ethylene oxide}]{\text{1.Mg/Et}_2\text{O}}$ Y $\xrightarrow{\text{KMnO}_4/\text{H}^{\oplus}}$ Z\\
(Palmitic acid) Choose correct option
\begin{qoptions}
\opt{A}{W is CH$_3$(CH$_2$)$_{14}$CH$_2$OH}
\opt{B}{X is CH$_3$(CH$_2$)$_{14}$CH$_2$Cl}
\opt{C}{Y is CH$_3$(CH$_2$)$_{16}$COOH}
\opt{D}{Z is CH$_3$(CH$_2$)$_{16}$CHO}
\end{qoptions}
\qn{5} \qfig{p2_chem_s1_q05_1}
(P) was molecular formula C$_7$H$_6$O$_2$ which forms `CO$_2$' gas with NaHCO$_3$. Choose the incorrect option(s)?
\begin{qoptions}
\opt{A}{Acidic strength order : (Q) $>$ (P) $>$ (R)}
\opt{B}{(P) has IUPAC name p-Nitrobenzoic acid}
\opt{C}{(R) has common name o-Bromonbenzoic acid}
\opt{D}{In formation of (Q) $\overset{\oplus}{\text{NO}_2}$ is electrophilic while for (R) Br$^{\oplus}$ is electrophite}
\end{qoptions}
\qn{6} In which of the following reaction, final major product gives intramolecular aldol condensation reaction ?
\begin{qoptions}
\opt{A}{\qfigi{p2_chem_s1_q06_A}}
\opt{B}{\qfigi{p2_chem_s1_q06_B}}
\opt{C}{\qfigi{p2_chem_s1_q06_C}}
\opt{D}{\qfigi{p2_chem_s1_q06_D}}
\end{qoptions}
\qn{7} \qfig{p2_chem_s1_q07_1}
Which of the following option is/are correct?
\begin{qoptions}
\opt{A}{P and Q are geometrical isomer}
\opt{B}{R and S are metamer}
\opt{C}{P to R is known as Beckmann rearrangment}
\opt{D}{X to P is known as Clemenson reduction}
\end{qoptions}
\qn{8} CH$_3-$C$\equiv$CH $\xrightarrow[\text{dil. H}_2\text{SO}_4]{\text{dil. HgSO}_4}$ (A) (Major)\\
Compound (A) and PhCHO can not be differentiated by :
\begin{qoptions}
\opt{A}{Fehling solution}
\opt{B}{Lucas test}
\opt{C}{2,4-DNP test}
\opt{D}{Iodoform test}
\end{qoptions}
\sect{SECTION-I(ii)}
\end{multicols}
\noindent\begin{minipage}{\textwidth}
\qn{9} Answer the following by appropriately matching the lists based on the information given in the paragraph\par\smallskip
\begin{center}\small\renewcommand{\arraystretch}{1.3}\setlength{\tabcolsep}{4pt}
\begin{tabular}{|l|c|l|p{6.2cm}|}
\hline & \textbf{List-I} & & \textbf{List-II} \\\hline
(I) & \qfigi{p2_chem_s2_q09_1}\ $\xrightarrow[\text{(iii) Ag}_2\text{O}/\Delta\ \text{(iv) HOH}]{\text{(i) SOCl}_2\ \text{(ii) CH}_2\text{N}_2}$ & (P) & Substrate is unable to evolve CO$_2$ on treatment with NaHCO$_3$ \\\hline
(II) & \qfigi{p2_chem_s2_q09_2} & (Q) & No dehydration but decarboxylation takes place in the reaction. \\\hline
(III) & \qfigi{p2_chem_s2_q09_3} & (R) & Racemic mixture obtained in the final product \\\hline
(IV) & \qfigi{p2_chem_s2_q09_4} & (S) & Number of C atom increases in the final product \\\hline
 &  & (T) & Reactant can undergo decarboxylation on heating \\\hline
 &  & (U) & Product obtained with lesser number of C atom as compare to reactant. \\\hline
\end{tabular}
\end{center}
\end{minipage}
\begin{multicols}{2}
Which of the following option contains all correct combinations?
\begin{qoptions}
\opt{A}{I $\to$ P; II $\to$ Q; III $\to$ T; IV $\to$ R}
\opt{B}{I $\to$ P; II $\to$ Q; III $\to$ R; IV $\to$ T}
\opt{C}{I $\to$ S; II $\to$ U; III $\to$ R; IV $\to$ P}
\opt{D}{I $\to$ S; II $\to$ U; III $\to$ P; IV $\to$ R}
\end{qoptions}
\end{multicols}
\noindent\begin{minipage}{\textwidth}
\qn{10} Answer the following by appropriately matching the lists based on the information given in the paragraph\par\smallskip
\begin{center}\small\renewcommand{\arraystretch}{1.3}\setlength{\tabcolsep}{4pt}
\begin{tabular}{|l|c|l|p{6.2cm}|}
\hline & \textbf{List-I} & & \textbf{List-II} \\\hline
(I) & \qfigi{p2_chem_s2_q10_1}\ $\xrightarrow[\text{(iii) Ag}_2\text{O}/\Delta\ \text{(iv) HOH}]{\text{(i) SOCl}_2\ \text{(ii) CH}_2\text{N}_2}$ & (P) & Substrate is unable to evolve CO$_2$ on treatment with NaHCO$_3$ \\\hline
(II) & \qfigi{p2_chem_s2_q10_2} & (Q) & No dehydration but decarboxylation takes place in the reaction. \\\hline
(III) & \qfigi{p2_chem_s2_q10_3} & (R) & Racemic mixture obtained in the final product \\\hline
(IV) & \qfigi{p2_chem_s2_q10_4} & (S) & Number of C atom increases in the final product \\\hline
 &  & (T) & Reactant can undergo decarboxylation on heating \\\hline
 &  & (U) & Product obtained with lesser number of C atom as compare to reactant. \\\hline
\end{tabular}
\end{center}
\end{minipage}
\begin{multicols}{2}
Which of the following is true. (about the substrate / reaction given above)?
\begin{qoptions}
\opt{A}{I $\to$ S; II $\to$ U; III $\to$ R; IV $\to$ P,R}
\opt{B}{I $\to$ R,S; II $\to$ Q,T; III $\to$ U; IV $\to$ P,U}
\opt{C}{I $\to$ S; II $\to$ Q,U; III $\to$ R; IV $\to$ P,S}
\opt{D}{I $\to$ R; II $\to$ T; III $\to$ U; IV $\to$ S}
\end{qoptions}
\end{multicols}
\noindent\begin{minipage}{\textwidth}
\qn{11} \par\smallskip
\begin{center}\small\renewcommand{\arraystretch}{1.3}\setlength{\tabcolsep}{4pt}
\begin{tabular}{|l|c|l|p{6.2cm}|}
\hline & \textbf{List-I} & & \textbf{List-II (A and B is differentiated by)} \\\hline
(I) & \qfigi{p2_chem_s2_q10_5} & (P) & Tollen's test \\\hline
(II) & \qfigi{p2_chem_s2_q10_6}\ \qfigi{p2_chem_s2_q10_7} & (Q) & Fehling test \\\hline
(III) & \qfigi{p2_chem_s2_q10_8} & (R) & NaHCO$_3$ \\\hline
(IV) & \qfigi{p2_chem_s2_q10_9} & (S) & Iodoform test \\\hline
 &  & (T) & FeCl$_3$ \\\hline
 &  & (U) & 2,4-DNP \\\hline
\end{tabular}
\end{center}
\end{minipage}
\begin{multicols}{2}
The incorrect option is :
\begin{qoptions}
\opt{A}{I $\to$ R; II $\to$ P; III $\to$ S; IV $\to$ S}
\opt{B}{I $\to$ R; II $\to$ Q; III $\to$ S; IV $\to$ S}
\opt{C}{I $\to$ R; II $\to$ S; III $\to$ S; IV $\to$ T}
\opt{D}{I $\to$ R; II $\to$ T; III $\to$ S; IV $\to$ U}
\end{qoptions}
\end{multicols}
\noindent\begin{minipage}{\textwidth}
\qn{12} \par\smallskip
\begin{center}\small\renewcommand{\arraystretch}{1.3}\setlength{\tabcolsep}{4pt}
\begin{tabular}{|l|c|l|p{6.2cm}|}
\hline & \textbf{List-I} & & \textbf{List-II (A and B is differentiated by)} \\\hline
(I) & \qfigi{p2_chem_s2_q11_1} & (P) & Tollen's test \\\hline
(II) & \qfigi{p2_chem_s2_q11_2}\ \qfigi{p2_chem_s2_q11_3} & (Q) & Fehling test \\\hline
(III) & \qfigi{p2_chem_s2_q11_4} & (R) & NaHCO$_3$ \\\hline
(IV) & \qfigi{p2_chem_s2_q11_5} & (S) & Iodoform test \\\hline
 &  & (T) & FeCl$_3$ \\\hline
 &  & (U) & 2,4-DNP \\\hline
\end{tabular}
\end{center}
\end{minipage}
\begin{multicols}{2}
The correct option is :
\begin{qoptions}
\opt{A}{I $\to$ R; II $\to$ S; III $\to$ S; IV $\to$ Q}
\opt{B}{I $\to$ R; II $\to$ P; III $\to$ S; IV $\to$ T}
\opt{C}{I $\to$ R; II $\to$ P; III $\to$ S; IV $\to$ Q}
\opt{D}{I $\to$ R; II $\to$ Q; III $\to$ P; IV $\to$ S}
\end{qoptions}
\sect{SECTION-II}
\qn{1} Number of reagents by which\par\smallskip\qfigi{p2_chem_s3_q01_1}\par\smallskip\qfigi{p2_chem_s3_q01_2}\par\smallskip acidification.

(i) [Aq.(NH$_3$)$_2$]OH \quad (ii) \qfigi{p2_chem_s3_q01_3}

(iii) K$_2$Cr$_2$O$_7$/H$^+$, $\Delta$ \quad (iv) Conc. $^-$OH

(v) KMnO$_4$, $\Delta$ \quad (vi) CuSO$_4$ + NH$_4$OH

(vii) Zn-Hg/HCl \quad (viii) Red P / HI

(ix) Alcoholic KCN \quad (x) \qfigi{p2_chem_s3_q01_4}
\qn{2} An organic compound (C$_8$H$_{10}$O$_2$) rotates plane-polarized light. It produces voilet color with neutral FeCl$_3$ solution. What is the total number of all the possible isomers (including stereoisomers) for this compound?
\qn{3} Scheme 1 and 2 describe the conversion of P to Q and R to S, respectively. Scheme 3 describes the synthesis of T from Q and S. The total number of Br atoms in a molecule of T is \underline{\hspace{1.5cm}}
\qfig{p2_chem_s3_q03_1}
\textbf{Scheme 2 :}
\qfig{p2_chem_s3_q03_2}
\qfig{p2_chem_s3_q03_3}
\qfullbegin
\qn{4} Identify total number of reactions incorrectly match with its product ?
\qfig{p2_chem_s3_q04_1}
\qfullsplit
(e) Me$-$C$\equiv$CH + Na/Liq NH$_3$ $\to$ Me$-$CH $=$ CH$_2$
\qfig{p2_chem_s3_q04_2}
\qfullend
\qn{5} How many of the following are correctly matched with their product?
\qfig{p2_chem_s3_q05_1}
\qfig{p2_chem_s3_q05_2}
\qfig{p2_chem_s3_q05_3}
\qfig{p2_chem_s3_q05_4}
\qfig{p2_chem_s3_q05_5}
\qfig{p2_chem_s3_q05_6}
\qn{6} \qfig{p2_chem_s3_q06_1}
Position of CHO group in A is $x$, no. of Br atoms in B is $y$ then $x+y=$ ?
\partsec{3}{MATHEMATICS}
\sect{SECTION-I(i)}
\qn{1} If $x\cos\alpha+y\sin\alpha=x\cos\beta+y\sin\beta=2a$ and $2\sin\Bigl(\dfrac\alpha2\Bigr)\sin\Bigl(\dfrac\beta2\Bigr)=1$ then
\begin{qoptions}
\opt{A}{$y^2=4a(a-x)$}
\opt{B}{$\cos\alpha+\cos\beta=\cos\alpha\cos\beta$}
\opt{C}{$\cos\alpha.\cos\beta=\dfrac{4a^2+y^2}{x^2+y^2}$}
\opt{D}{$(\cos\alpha+\cos\beta)=\dfrac{4ax}{x^2+y^2}$}
\end{qoptions}
\qn{2} If $\cos\beta$ is geometric mean between $\sin\alpha$ and $\cos\alpha$, where $0<\alpha,\beta<\dfrac\pi2$, then $\cos2\beta=$
\begin{qoptions}
\opt{A}{$-2\sin^2\Bigl(\dfrac\pi4-\alpha\Bigr)$}
\opt{B}{$-2\cos^2\Bigl(\dfrac\pi4+\alpha\Bigr)$}
\opt{C}{$2\sin^2\Bigl(\dfrac\pi4+\alpha\Bigr)$}
\opt{D}{$2\cos^2\Bigl(\dfrac\pi4-\alpha\Bigr)$}
\end{qoptions}
\qn{3} If $b>a>0$, then the expression $E$ given by $E=\dfrac{1}{\sqrt{b-a}}\cdot\dfrac{\sqrt{\frac{b-a}{a}}\sin x}{\sqrt{\Bigl[1+\Bigl(\sqrt{\frac{b-a}{a}}\sin x\Bigr)^2\Bigr]}}\cdot\sqrt{a+b\tan^2x}$ must be equal to
\begin{qoptions}
\opt{A}{$\tan x$}
\opt{B}{$|\tan x|$}
\opt{C}{$\dfrac{\sin x}{|\cos x|}$}
\opt{D}{unity if $x=\dfrac{11\pi}{4}$}
\end{qoptions}
\qn{4} If $\sin(x+20^\circ)=2\sin x\cos40^\circ$, where $x\in(0,\pi/2)$, then which of the following hold(s) good
\begin{qoptions}
\opt{A}{$\cos2x=\dfrac12$}
\opt{B}{$\operatorname{cosec}4x=2$}
\opt{C}{$\sec\dfrac x2=\sqrt6-\sqrt2$}
\opt{D}{$\tan\dfrac x2=\bigl(2-\sqrt3\bigr)$}
\end{qoptions}
\qn{5} If $\cos\alpha+\cos\beta=a$,\\
$2(\cos^2\alpha+\cos^2\beta)=b+2$,\\
$\cos3\alpha+\cos3\beta=c$, then which of the following options is/are correct ?
\begin{qoptions}
\opt{A}{$\cos\alpha\cos\beta=\dfrac{a^2}{2}-\Bigl(\dfrac{b+2}{4}\Bigr)$}
\opt{B}{$2a^3+c=3a(1+b)$}
\opt{C}{$a+b+c=3abc$}
\opt{D}{$2(1+b)=\dfrac{\sin3\alpha}{\sin\alpha}+\dfrac{\sin3\beta}{\sin\beta}$, for $\alpha,\beta\neq n\pi$, $n\in I$}
\end{qoptions}
\qn{6} A given line $L_1$ intersects $x$ and $y$ axes at $P$ and $Q$ respectively. A variable line $L_2$, perpendicular to $L_1$, cuts $x$ and $y$ axes at $R$ and $S$ respectively. The locus of the point of intersection of the lines $PS$ and $QR$ is a curve passing through the
\begin{qoptions}
\opt{A}{Origin}
\opt{B}{Point $P$}
\opt{C}{Point $Q$}
\opt{D}{Point $R$}
\end{qoptions}
\qn{7} If $A(3,4)$ and $B(-5,-2)$ are the extremities of the base of an isosceles triangle $ABC$ with $\tan C=2$, then point $C$ can be
\begin{qoptions}
\opt{A}{$\Bigl(\dfrac{3\sqrt5-1}{2},-\bigl(1+2\sqrt5\bigr)\Bigr)$}
\opt{B}{$\Bigl(-\dfrac{\bigl(3\sqrt5+5\bigr)}{2},3+2\sqrt5\Bigr)$}
\opt{C}{$\Bigl(\dfrac{\bigl(3\sqrt5+5\bigr)}{2},3-2\sqrt5\Bigr)$}
\opt{D}{$\Bigl(\dfrac{\bigl(3\sqrt5+1\bigr)}{2},-\bigl(1+2\sqrt5\bigr)\Bigr)$}
\end{qoptions}
\qn{8} Two fixed points $P$ and $Q$ are 4 units apart and are on same side of a variable line $L$. Let $PM$ and $QN$ are perpendicular distance of $P$ and $Q$ from line $L$ satisfying equation $PM+3QN=4$, then line $L$ always touches a circle $C$ then which among the following option(s) is/are correct for $C$ ?
\begin{qoptions}
\opt{A}{the centre of circle $C$ lies on line $PQ$}
\opt{B}{the radius of circle is 1}
\opt{C}{the radius of circle is 2}
\opt{D}{the center of the circle $C$ lies on perpendicular bisector of $PQ$}
\end{qoptions}
\sect{SECTION-I(ii)}
\end{multicols}
\noindent\begin{minipage}{\textwidth}
\qn{9} Let $A(x_1,y_1)$, $B(x_2,y_2)$, $C(x_3,y_3)$ be three distinct points lying on circle $S:\ x^2+y^2=1$ such that $x_1x_2+y_1y_2+x_2x_3+y_2y_3+x_3x_1+y_3y_1=-\dfrac32$\par\smallskip
\begin{center}\small\renewcommand{\arraystretch}{1.25}\setlength{\tabcolsep}{4pt}
\begin{tabular}{|l|p{10.5cm}|l|c|}
\hline & \textbf{List-I} & & \textbf{List-II} \\\hline
(I) & Let $P$ be any arbitrary point lying on $S$, then $(PA)^2+(PB)^2+(PC)^2=$ & (P) & 3 \\\hline
(II) & Let the perpendicular dropped from point `A' to BC meets $S$ at $Q$ and $\angle OBQ=\dfrac{\pi}{k}$, where `O' is origin, then $k=$ & (Q) & 4 \\\hline
(III) & Let $R$ be the point lying on line $x+y=2$, at the minimum distance from $S$ and the square of maximum distance of $R$ from $S$ is $a+b\sqrt b$, then $a+b=\ldots..$ (Given $a$ and $b$ are distinct natural number) & (R) & 6 \\\hline
(IV) & Let $I$ and $G$ represent incentre and centroid of $\Delta ABC$ respectively, then $IA+IB+IC+GA+GB+GC=$ & (S) & 5 \\\hline
 &  & (T) & 2 \\\hline
 &  & (U) & 7 \\\hline
\end{tabular}
\end{center}
\end{minipage}
\begin{multicols}{2}
Which of the following is the only correct combination :
\begin{qoptions}
\opt{A}{I $\to$ R; II $\to$ P; III $\to$ S}
\opt{B}{I $\to$ R; II $\to$ P; IV $\to$ T}
\opt{C}{I $\to$ R; III $\to$ U; IV $\to$ R}
\opt{D}{II $\to$ P; III $\to$ T; IV $\to$ R}
\end{qoptions}
\end{multicols}
\noindent\begin{minipage}{\textwidth}
\qn{10} Let $A(x_1,y_1)$, $B(x_2,y_2)$, $C(x_3,y_3)$ be three distinct points lying on circle $S:\ x^2+y^2=1$ such that $x_1x_2+y_1y_2+x_2x_3+y_2y_3+x_3x_1+y_3y_1=-\dfrac32$\par\smallskip
\begin{center}\small\renewcommand{\arraystretch}{1.25}\setlength{\tabcolsep}{4pt}
\begin{tabular}{|l|p{10.5cm}|l|c|}
\hline & \textbf{List-I} & & \textbf{List-II} \\\hline
(I) & Let $P$ be any arbitrary point lying on $S$, then $(PA)^2+(PB)^2+(PC)^2=$ & (P) & 3 \\\hline
(II) & Let the perpendicular dropped from point `A' to BC meets $S$ at $Q$ and $\angle OBQ=\dfrac{\pi}{k}$, where `O' is origin, then $k=$ & (Q) & 4 \\\hline
(III) & Let $R$ be the point lying on line $x+y=2$, at the minimum distance from $S$ and the square of maximum distance of $R$ from $S$ is $a+b\sqrt b$, then $a+b=\ldots..$ (Given $a$ and $b$ are distinct natural number) & (R) & 6 \\\hline
(IV) & Let $I$ and $G$ represent incentre and centroid of $\Delta ABC$ respectively, then $IA+IB+IC+GA+GB+GC=$ & (S) & 5 \\\hline
 &  & (T) & 2 \\\hline
 &  & (U) & 7 \\\hline
\end{tabular}
\end{center}
\end{minipage}
\begin{multicols}{2}
Which of the following is the only incorrect combination :
\begin{qoptions}
\opt{A}{I $\to$ R; II $\to$ P}
\opt{B}{III $\to$ S; IV $\to$ R}
\opt{C}{I $\to$ R; III $\to$ S}
\opt{D}{II $\to$ R; IV $\to$ T}
\end{qoptions}
\end{multicols}
\noindent\begin{minipage}{\textwidth}
\qn{11} Match the set of values of $x$ in List-II which satisfy the equations in List-I.\\ Let $D(0,\sqrt3)$, $E(1,0)$, $F(-1,0)$ be the feet of perpendiculars dropped from vertices to the opposite sides $BC$, $CA$, $AB$ respectively of triangle $ABC$.\par\smallskip
\begin{center}\small\renewcommand{\arraystretch}{1.25}\setlength{\tabcolsep}{4pt}
\begin{tabular}{|l|p{10.5cm}|l|c|}
\hline & \textbf{List-I} & & \textbf{List-II} \\\hline
(I) & The ratio of the inradius of $\Delta ABC$ to the inradius of $\Delta DEF$ is & (P) & 2 \\\hline
(II) & Let `H' be the orthocentre of $\Delta ABC$, then the greatest integer which is less than or equal to square of the length $AH$ is & (Q) & 3 \\\hline
(III) & The square of the sum of ordinates of points $A$, $B$ and $C$ is & (R) & 4 \\\hline
(IV) & The length of side $AB$ of $\Delta ABC$ is & (S) & 5 \\\hline
 &  & (T) & 1 \\\hline
 &  & (U) & 6 \\\hline
\end{tabular}
\end{center}
\end{minipage}
\begin{multicols}{2}
Which of the following is the only correct combination :
\begin{qoptions}
\opt{A}{I $\to$ P; II $\to$ S}
\opt{B}{II $\to$ S; III $\to$ T}
\opt{C}{III $\to$ T; IV $\to$ R}
\opt{D}{I $\to$ P; IV $\to$ U}
\end{qoptions}
\end{multicols}
\noindent\begin{minipage}{\textwidth}
\qn{12} Match the set of values of $x$ in List-II which satisfy the equations in List-I.\\ Let $D(0,\sqrt3)$, $E(1,0)$, $F(-1,0)$ be the feet of perpendiculars dropped from vertices to the opposite sides $BC$, $CA$, $AB$ respectively of triangle $ABC$.\par\smallskip
\begin{center}\small\renewcommand{\arraystretch}{1.25}\setlength{\tabcolsep}{4pt}
\begin{tabular}{|l|p{10.5cm}|l|c|}
\hline & \textbf{List-I} & & \textbf{List-II} \\\hline
(I) & The ratio of the inradius of $\Delta ABC$ to the inradius of $\Delta DEF$ is & (P) & 2 \\\hline
(II) & Let `H' be the orthocentre of $\Delta ABC$, then the greatest integer which is less than or equal to square of the length $AH$ is & (Q) & 3 \\\hline
(III) & The square of the sum of ordinates of points $A$, $B$ and $C$ is & (R) & 4 \\\hline
(IV) & The length of side $AB$ of $\Delta ABC$ is & (S) & 5 \\\hline
 &  & (T) & 1 \\\hline
 &  & (U) & 6 \\\hline
\end{tabular}
\end{center}
\end{minipage}
\begin{multicols}{2}
Which of the following is the only incorrect combination :
\begin{qoptions}
\opt{A}{I $\to$ P; II $\to$ S}
\opt{B}{III $\to$ Q; IV $\to$ R}
\opt{C}{I $\to$ S; III $\to$ P}
\opt{D}{II $\to$ S; IV $\to$ R}
\end{qoptions}
\sect{SECTION-II}
\qn{1} Let $x$ be in radians with $0<x<\dfrac\pi2$. If $\sin(2\sin x)=\cos(2\cos x)$ ; then the value of $(\tan x+\cot x)$ can be written as $\pm\Bigl(\dfrac{a}{\pi^c-b}\Bigr)$ where $a,b,c\in N$. Then the value of $\Bigl(\dfrac{a+b+c}{25}\Bigr)$ is
\qn{2} $AB$ is a line segment of length 48cm and $C$ is its mid-point. If three semicircles are drawn at $AB$, $AC$ and $CB$ using as diameters, then radius of the circle inscribed in the space enclosed by three semicircles is
\qn{3} If $P(x_0,y_0)$ be the coordinates of points on curve $5x^2-6xy+5y^2-4=0$. Which is nearest to origin and $d(x_0,y_0)$ denotes its distance from origin then the value of $\dfrac{|x_0|+|y_0|}{\bigl(d(x_0,y_0)\bigr)^2}$ is equal to
\qn{4} Let $L_1$, $L_2$ and $L_3$ be the length of tangents drawn from a point $P$ to the circles $x^2+y^2=4$, $x^2+y^2-4x=0$ and $x^2+y^2-4y=0$ respectively. If $L_1^4=L_2^2L_3^2+16$, then locus of $P$ are the curve $C_1$ (a straight line) and $C_2$ (a circle). Let the circumcentre of the triangle formed by $C_1$ and other two lines which are inclined at an angle of $45^\circ$ with $C_1$ and are tangents to $C_2$ is $(\alpha,\beta)$ then $(\alpha+2\beta)$ equals
\qn{5} Let $r_1$ and $r_2$ be the radii of the largest and smallest circles, respectively, which pass through the point $(-4,1)$ and having their centres on the circumference of the circle $x^2+y^2+2x+4y-4=0$. If $\dfrac{r_1}{r_2}=a+b\sqrt2$ then $a+b$ is equal to :
\qn{6} A circle passes through the vertex $C$ of a rectangle $ABCD$ and touches its sides $AB$ and $AD$ at $M$ and $N$, respectively. If the distance from $C$ to the line-segment $MN$ is equal to 5 units, then $\dfrac15$ (area of the rectangle $ABCD$) =

\end{multicols}
\end{document}
